% The prolate mechanism is sign-blind --- a sharpening of the positivity obstruction.
%
% DRAFT.  Assembled 2026-08-12 (work item mg-8b7b) from the notes in ../notes/.
% Not a submission.  Publication is not decided by this document.
%
% Build:  pdflatex positivity-obstruction.tex   (twice, for cross-references)

\documentclass[11pt,a4paper]{article}

\usepackage[margin=1.05in]{geometry}
\usepackage{amsmath,amssymb,amsthm}
\usepackage[T1]{fontenc}
\usepackage{microtype}
\usepackage{url}
\usepackage{xcolor}
\usepackage[colorlinks=true,linkcolor=blue!45!black,citecolor=blue!45!black,%
            urlcolor=blue!45!black]{hyperref}
\sloppy

\theoremstyle{plain}
\newtheorem{theorem}{Theorem}[section]
\newtheorem{proposition}[theorem]{Proposition}
\newtheorem{lemma}[theorem]{Lemma}
\newtheorem{corollary}[theorem]{Corollary}
\theoremstyle{definition}
\newtheorem{observation}[theorem]{Observation}
\newtheorem{hypothesis}{Hypothesis}
\renewcommand{\thehypothesis}{H\arabic{hypothesis}}
\theoremstyle{remark}
\newtheorem{remark}[theorem]{Remark}

% --- status and attribution tags ------------------------------------------
% Every load-bearing statement in this paper carries one status tag and one
% attribution tag.  See section 1.5 for what the words mean.
\newcommand{\st}[1]{{\normalfont\sffamily\footnotesize[#1]}}
\newcommand{\proved}{\st{proved}}
\newcommand{\announced}{\st{announced}}
\newcommand{\observed}{\st{observed}}
\newcommand{\conj}{\st{conjectural}}
\newcommand{\ours}{{\normalfont\sffamily\footnotesize[ours]}}
\newcommand{\theirs}[1]{{\normalfont\sffamily\footnotesize[#1]}}

\newcommand{\R}{\mathbb{R}}
\newcommand{\Z}{\mathbb{Z}}
\newcommand{\Q}{\mathbb{Q}}
\newcommand{\A}{\mathbb{A}}
\newcommand{\E}{\mathcal{E}}
\newcommand{\F}{\mathcal{F}}
\newcommand{\PS}{\mathit{PS}}
\newcommand{\Sonin}{\mathbf{S}}
\newcommand{\QW}{QW}
\newcommand{\file}[1]{\path{#1}}

% keep macros out of PDF bookmark strings
\pdfstringdefDisableCommands{%
  \def\QW{QW}\def\E{E}\def\F{F}\def\R{R}\def\Q{Q}\def\Z{Z}\def\A{A}%
  \def\PS{PS}\def\Sonin{S}\def\file#1{#1}%
}

\title{The prolate mechanism is sign-blind:\\
a sharpening of the positivity obstruction}

\author{\normalsize (draft --- author line deliberately left blank)}
\date{Draft of 12 August 2026}

\begin{document}
\maketitle

\begin{abstract}
\noindent
Connes and Consani exhibit, numerically, a family of minuscule eigenvalues
$s(\mu)$ of the semi-local Weil quadratic form $\QW_\lambda$ restricted to test
functions supported in $[\lambda^{-1},\lambda]$, $\mu=\lambda^2$, and locate the
mechanism behind them: the radical of $\QW_\lambda$ contains the range of the
summation map $\E$, and the failure of the double support condition is a
Slepian concentration defect.  This paper does three things.  First, it
assembles what is \emph{proved} in that circle of ideas and separates it from
what is announced, observed, or conjectural, claim by claim and number by
number.  Second, it makes the size of $s(\mu)$ quantitative: $s(\mu)$ is
governed by the prolate concentration defect at index~$4$ --- not index~$0$ ---
with $s/(1-\chi_2)$ between $7$ and $13$ over $\mu=5,\dots,12$, and the
underlying construction is carried to an identity, $\QW_\lambda$ at the prolate
vector \emph{is} the Weil form of the part of $\E(\phi)$ that falls outside the
window.  Third, and this is the point, it observes that the entire mechanism is
invariant under $\QW_\lambda\mapsto-\QW_\lambda$.  A kernel is a kernel either
way; a concentration defect is a magnitude.  So the mechanism explains why
$|s(\mu)|$ is minuscule and says nothing whatever about why $s(\mu)>0$ --- and
that sign-blindness is exactly what forces the missing lower bound to carry the
whole arithmetic content.  Since a lower bound on $s$ of any shape is equivalent
to the Riemann Hypothesis, the decay rate $4\pi$ cannot be made an unconditional
theorem, and no computation can change that.  What remains is a genuine
sharpening of the obstruction, stated as three explicit constraints on method:
no fixed-relative-error argument can work; no prime below $\mu$ may be dropped
and no tail estimated at any cutoff; and the difficulty is cancellation, not
magnitude --- the archimedean deficit is bounded by
$\log\pi-\psi(\tfrac14)=5.3721834\ldots$ for every $\mu$.
\end{abstract}

\noindent
\textbf{Status of this document.}  This is a draft.  It is not a submission, and
nothing in it has been sent anywhere.  It assembles the working notes produced
in this project's repository on 12 August 2026 (listed in \S\ref{sec:provenance});
each section names the note it comes from, so every claim below can be traced to a source or to a script
without asking the authors anything.  Where a claim is not ours, it is marked so
at the point of use.  The five load-bearing numbers of
\S\ref{sec:deficit}--\S\ref{sec:scale} --- the archimedean bound and its
saturation, the deficit against the repair, the reversed signs on the
near-radical, the fitted rate, and the index-4 ratio table --- have since been
recomputed by a second implementation of ours, written from the definitions
without reading the first and sharing with it no arithmetic library, no
transcendental function and no eigensolver; all five reproduce to every printed
digit.  What that does and --- as sharply --- does
\emph{not} establish is set out in \S\ref{sec:recheck}: both implementations
read the same paper, so a shared misreading of it is invisible to the exercise.

\tableofcontents

% ==========================================================================
\section{Introduction}
% ==========================================================================

\subsection{The object}

Let $\lambda>1$, $\mu=\lambda^2$, $L=2\log\lambda=\log\mu$.  Connes and
Consani's \emph{semi-local Weil quadratic form}
(\cite{CC2023}, source line \file{Spectraltriples.tex:169}, eq.\ \texttt{weilQ}) is
\begin{equation}\label{eq:QW}
  \QW(f,g)=\sum_{\frac12+is\in Z}\overline{\widehat f(\bar s)}\,\widehat g(s),
  \qquad
  \widehat f(s)=\F_\mu(f)(s)=\int_{\R_+^*}f(u)\,u^{-is}\,d^*u ,
\end{equation}
with $Z$ the set of non-trivial zeros of $\zeta$; $\QW_\lambda$ denotes its
restriction to test functions supported in $[\lambda^{-1},\lambda]$.  The symbol
$\QW_\lambda$ is \emph{one symbol} --- $Q$ for quadratic form, $W$ for Weil.
There is no operator $Q$ in this literature (\S\ref{sec:onesymbol}).  By Weil's
criterion, positivity of $\QW_\lambda$ for all $\lambda>1$ is equivalent to the
Riemann Hypothesis (\cite{CC2023} at \file{:169}; \cite{Connes2026} at
\file{rhready.tex:1145}). \proved

By Proposition \texttt{Hilbert} of \cite{CC2023} (\file{:289}, proved there),
$\QW_\lambda$ has the Hilbert-space form
\begin{equation}\label{eq:hilbert}
  \QW_\lambda(f,f)=
  \underbrace{\int|\widehat f(t)|^2\,\frac{2\vartheta'(t)}{2\pi}\,dt
  \;+\;2\,\Re\!\big(\widehat f(\tfrac i2)\overline{\widehat f(-\tfrac i2)}\big)}
  _{\textstyle \sigma^{\mathrm{arch}}}
  \;-\;\sum_{1<n\le\mu}\Lambda(n)\,\langle f\mid V(n)f\rangle ,
\end{equation}
$\vartheta$ the Riemann--Siegel theta function and $V(n)$ the compression to
$[\lambda^{-1},\lambda]$ of $n^{-1/2}(\vartheta_n+\vartheta_{n^{-1}})$.  In the
even sector (parity under $u\mapsto u^{-1}$) the truncated form is a matrix
$\sigma^+$, and we write
\[
  s(\mu):=\lambda_{\min}(\sigma^+),
\]
which is Connes--Consani's $s(L)$ at \file{:645} and, read as
$\min(s^+,s^-)$, Connes' $\epsilon(\lambda)$ at
\file{rhready.tex:1145}--\file{:1147}.

\subsection{The question, and the thesis}

Connes and Consani report that $s(\mu)$ decays exponentially in $\mu$ and
display it as a figure; at $\mu=11$ the smallest positive eigenvalue is
$2.389\times10^{-48}$ (\file{:178}).  Connes' 2026 survey reports ``a striking
similarity'' between $\epsilon(\lambda)$ and the angular function
$1-\chi_2(\lambda)$, again as a figure, with no constant
(\file{rhready.tex:1149}). \observed

The natural programme is to convert that similarity into an estimate and the
estimate into a proof of positivity.  This paper reports how far that goes and
where it stops, and the stopping point is structural rather than technical:

\begin{quote}
\textbf{Thesis.}  The prolate mechanism --- the radical of $\QW_\lambda$
contains $\operatorname{ran}\E$, the near-radical is the failure of a double
support condition, and the size of the residue is a concentration defect --- is
\emph{entirely invariant} under $\QW_\lambda\mapsto-\QW_\lambda$.  A kernel is a
kernel either way; a concentration defect is a magnitude.  So it explains why
$|s(\mu)|$ is minuscule and says nothing whatever about why $s(\mu)>0$.  That
sign-blindness is what forces the missing lower bound to carry the whole
arithmetic content, and the arithmetic content is~RH.
\end{quote}

Everything below is either evidence for that claim or a consequence of it.  It
is not an apology for a failed attempt: the sign-blindness is a property of the
mechanism that can be checked one statement at a time, and \S\ref{sec:houserule}
checks it on every principal claim of this paper, including the ones we are
pleased with.

\subsection{Results, and whose they are}
\label{sec:whose}

The following is the whole of what this paper asserts, with attribution.  Rows
marked \emph{ours} are established in the project notes named in
\S\ref{sec:provenance}; the rest are in the literature and are recalled here
because the argument needs them.

\begin{center}\small
\begin{tabular}{@{}p{0.50\textwidth}p{0.20\textwidth}p{0.22\textwidth}@{}}
\hline
statement & status & whose \\
\hline
$W_\infty(g\star g^*)\ge\operatorname{Tr}(\vartheta(g)\Sonin\vartheta(g)^*)$ for
support in $[2^{-1/2},2^{1/2}]$; the sign is a square
& proved & Connes--Consani \cite{CC2021} \\
$V(n)$ is unitarily equivalent to $-V(n)$; its spectrum is exactly symmetric
about $0$, with $\|V(n)\|=2n^{-1/2}\cos(\pi/(m+2))$
& proved & \emph{ours} (\S\ref{sec:indefinite}) \\
At a finite place the two cutoff projections commute over $\Q_p$: no plunge, no
prolate family
& proved & Burnol \cite{Burnol1999}; reading \emph{ours} \\
Semilocally the angle operator has no spectral description, because $P^S_T$ cuts
off the product module and does not factor
& proved (definitional) & \emph{ours}, from \cite{Connes2026,CCM2023} \\
Archimedean-only positivity fails past $\mu=2.2710$
& observed, reproduced & Connes--Consani \cite{CC2023}; reproduction \emph{ours} \\
An explicit truncation-free test function on which the archimedean form is
negative for $\mu\ge3.5581$, and the uncertainty-principle reason
& proved (for that function) & \emph{ours} (\S\ref{sec:witness}) \\
$D(\mu)=-\lambda_{\min}(\sigma^{\mathrm{arch}})<\log\pi-\psi(\tfrac14)=5.3721834\ldots$
for every $\mu$, and it saturates
& proved & \emph{ours} (\S\ref{sec:deficit}) \\
On the near-radical direction the signs reverse: the archimedean half is the
credit and the primes are the debit
& observed & \emph{ours} (\S\ref{sec:thirddirection}) \\
$1-\chi_2\sim\frac{2^{14}\sqrt2\,\pi^5}{3}\mu^{9/2}e^{-4\pi\mu}$, and
$\epsilon(\lambda)$ behaves like $1-\chi_2$
& proved (asymptotic); observed (similarity) & Connes \cite{Connes2026}, from
Fuchs \cite{Fuchs1964} \\
$s(\mu)$ is governed by the defect at prolate index~$4$, not $0$;
$s/(1-\chi_2)\in[7.1,\,9.9]$ after extrapolation, with no trend
& observed (arbitrary precision) & \emph{ours} (\S\ref{sec:confront}) \\
The rate $4\pi$ \emph{cannot} discriminate the index; the power $9/2$ and the
constant can, and both are published
& proved & \emph{ours} as an observation (\S\ref{sec:index}) \\
$\QW_\lambda(g,g)=\sum_{Z}|\F_\mu r(s)|^2$ with $r$ the spill outside the window
& proved & \emph{ours}, from \cite{CC2023} (\S\ref{sec:identity}) \\
Under \eqref{H0}+\eqref{H1}, $\limsup\mu^{-1}\log s(\mu)\le-4\pi$
& conditional & \emph{ours} (\S\ref{sec:identity}) \\
A lower bound $s(\mu)\ge\kappa(\mu)>0$ for all $\mu$ implies RH; the two-sided
rate is therefore unreachable
& proved & \emph{ours} as an observation, from Weil and \cite{Connes2026} \\
Three constraints on method (\S\ref{sec:sharpening})
& proved / observed, per constraint & \emph{ours} \\
The five load-bearing numbers of \S\ref{sec:deficit}--\S\ref{sec:scale}
reproduce under a second implementation sharing no arithmetic library, no
transcendental function and no eigensolver --- but the same reading of
\cite{CC2023}
& verified; not the shared reading
& \emph{ours} (\S\ref{sec:recheck}) \\
The printed closed form (\texttt{h02ev}) at \file{Spectraltriples.tex:474} is a
factor $2$ from the table at \file{:444} it is derived from; no result here
depends on it
& observed & \emph{ours} (\S\ref{sec:erratum}), on \cite{CC2023} \\
\hline
\end{tabular}
\end{center}

Two attributions are worth stating separately because it would be easy to appear
to claim them.  \emph{The identification of the controlling scale is not ours}:
Connes names $1-\chi_2$ and gives its asymptotic in the February 2026 survey.
\emph{The construction of the near-radical vectors is not ours}: the two-mode
prolate combination $\phi_2=\psi_2\psi_0(0)-\psi_0\psi_2(0)$ is
Connes--Consani's, at \file{Spectraltriples.tex:744}.  What is ours in
\S\ref{sec:scale} is the quantitative confrontation, the index argument, the
identity, and the observation that the rate is a non-discriminating parameter.

\subsection{Routes closed, and at what threshold}
\label{sec:closed}

A reader who knows the programme will want, in one place, which methods are
excluded and by what.  Each row is argued in the section named.

\begin{center}\small
\begin{tabular}{@{}p{0.30\textwidth}p{0.34\textwidth}p{0.26\textwidth}@{}}
\hline
route & why it is closed & threshold / status \\
\hline
Assemble positivity place by place
& the archimedean term alone goes negative; no $-\Lambda(n)V(n)$ is of one sign
& $\mu>2.2710$ \observed; explicit witness at $\mu\ge3.5581$ \proved{};
  $V(n)$ indefinite at every $\mu$ \proved{} (\S\ref{sec:noplacewise}) \\
Get the sign from a compression of a projection, as at the archimedean place
& a compression $\vartheta(g)\Pi\vartheta(g)^*$ is positive; $V(n)$'s spectrum is
exactly symmetric about $0$
& all $\mu$, all $n$ \proved{} (\S\ref{sec:indefinite}) \\
Bound either side to a fixed relative error, or to a fixed power of $\mu$
& the two sides agree to $\approx5.46\mu$ decimal digits; the required relative
accuracy is $e^{-4\pi\mu}/(0.025\log\mu)$
& fails at
  $\mu\gtrsim(k\log\mu+\log\varepsilon^{-1})/5.4575$ \observed{}
  (\S\ref{sec:sharpening}) \\
Drop a prime, or estimate the tail $\sum_{p>P}$ crudely
& omitting any single $p^m<\mu$ makes the form negative by $0.10$ to $0.84$, i.e.\
$53$ orders above the residue, and the cost is not ordered by size
& every cutoff $P<\mu$ \observed{} (\S\ref{sec:loadbearing}) \\
Argue that the archimedean side is the thing that runs away
& it does not: $D(\mu)<5.3721834\ldots$ for every $\mu$
& all $\mu$ \proved{} (\S\ref{sec:deficit}) \\
Look for a plunge / prolate family at a single finite place
& over $\Q_p$ the two cutoff projections commute; the angle operator is trivial
& all $\Lambda=p^n$ \proved{} (\S\ref{sec:burnol}) \\
Prove the decay rate $4\pi$ as a two-sided unconditional theorem
& its lower half is a positive lower bound on $\QW_\lambda$ for all $\lambda$,
which is equivalent to RH
& all $\mu$ \proved{} (\S\ref{sec:boundary}) \\
Settle the rate by a wider grid or a converged truncation
& a computation cannot supply a lower bound at a $\mu$ it has not computed, and
the statement quantifies over all $\mu$
& permanent \proved{} (\S\ref{sec:boundary}) \\
\hline
\end{tabular}
\end{center}

\subsection{Standards of evidence used throughout}
\label{sec:standards}

These are not decoration; the corrections recorded in \S\ref{sec:provenance}
were each caught by one of them.

\paragraph{Status words.}  Every load-bearing sentence carries one.
\emph{Proved}: a theorem with a proof in a published or arXived paper, cited by
label and source line, or proved here.  \emph{Announced}: the authors say it can
be done, or defer it to a forthcoming paper; no proof given.  \emph{Observed}:
supported by computation, presented as evidence, and not a theorem.
\emph{Conjectural}: neither.  The distinction is not pedantry: an earlier draft
of this project's notes read a stated strategy as a companion theorem, and the
same class of error was then made three more times.

\paragraph{Reproduction versus verification.}  A number we recomputed
independently at a precision where the computation is not delicate is a
\emph{verification}.  A number we recomputed in a regime where the arithmetic is
tight, and which agrees with a published one, is a \emph{reproduction} --- it
confirms that we and the source are computing the same thing, and does not
certify the digits.  Both appear below and each is labelled.  For example, the
sign change at $\mu=2.2710$ is a verification (entries $O(1)$, spectrum down to
$1.3\times10^{-3}$); the match with $2.389\times10^{-48}$ at $\mu=11$ is a
reproduction whose \emph{exponent} is settled and whose mantissa is not
converged (\S\ref{sec:theirnumber}).

\paragraph{Corroboration, and what it cannot reach.}  A number recomputed by the
same code is not checked at all, and a number recomputed by a second
implementation is checked only against the causes of error the two do not share.
So a corroboration is worth exactly the list of what is \emph{not} shared, and it
is worthless against anything that is.  \S\ref{sec:recheck} gives both lists for
our own recheck of \S\ref{sec:deficit}--\S\ref{sec:scale}, in that order and at
the same length.

\paragraph{Second-hand numbers.}  A number that reaches us through a summary
rather than through the source is either re-opened at the source or marked.  One
class of external fact is used here without being read: Fuchs' 1964 asymptotic
\cite{Fuchs1964}, taken in the form Connes cites it and then checked against an
independent arbitrary-precision computation of $\Lambda_0,\Lambda_2,\Lambda_4$
(\S\ref{sec:fuchs}).  Yoshida's 1992 paper \cite{Yoshida1992} is likewise not
read; it is the one place a sharpness statement for the archimedean theorem
could be hiding (\S\ref{sec:notheorem}).

\paragraph{The standing test (``the house rule'').}  \emph{Is the statement
false for $-\QW_\lambda$?}  If not, it is sign-blind, and it cannot be evidence
about the direction of an inequality however sharp it is.  This test is applied
to the present paper in \S\ref{sec:houserule}, and it is the source of the
thesis in \S1.2.

% ==========================================================================
\section{The objects, and the literature they belong to}
\label{sec:objects}
% ==========================================================================

This section claims nothing.  It exists so that a reader can see which objects
below are Connes--Consani's --- almost all of them --- and so that the small
residue that is not theirs can be identified precisely.

\subsection{\texorpdfstring{$\QW_\lambda$}{QW-lambda} is one symbol}
\label{sec:onesymbol}

In \cite{CC2023}, $\mathbf{W}_\lambda$ (in bold, four occurrences) is the
\emph{prolate differential operator}
$(\mathbf{W}_\lambda\psi)(q)=-\partial((\lambda^2-q^2)\partial)\psi+(2\pi\lambda q)^2\psi$
of eq.\ (3.3), self-adjoint, positive, with the prolate spheroidal wave
functions as eigenfunctions; while $W_\lambda$ occurs $41$ times, \emph{every
one of them inside} $\QW_\lambda$, which is defined at \file{:169} as ``the
restriction $\QW_\lambda$ of the sesquilinear form\ldots''.  Two different
objects in one section, separated by a font. \proved{} (read in the arXiv \LaTeX{}
source.)

The consequence used below: there is no projection $Q$ to be applied after a
form.  The two conditions $f(0)=\widehat f(0)=0$ that one might expect such a
projection to impose are already the definition of the domain
$\mathcal S_0^{ev}$ of the summation map $\E$, and they are conditions on the
test function imposed \emph{before} $\E$ is applied.

\subsection{Object by object}

With $\E(f)(x)=x^{1/2}\sum_{n>0}f(nx)$ on
$\mathcal S_0^{ev}=\{f\ \text{even Schwartz}:f(0)=\widehat f(0)=0\}$, the
following are Connes--Consani's objects in Connes--Consani's letters
(\cite{CC2023} \S3): the map $\E$ itself (eq.\ (3.1)); the projections
$P_\lambda$, $\widehat P_\lambda$ and their angle operator
$P_\lambda\widehat P_\lambda P_\lambda$, diagonalised by the prolates; the
characteristic values $\chi(\mu,m)$ with
$\F_{e_\R}\psi_{m,\lambda}=\chi_m\psi_{m,\lambda}$ on $[-\lambda,\lambda]$ and
$\chi_m\simeq(-1)^m$; the count $\nu(\mu)\sim2\mu$ of small eigenvalues of the
angle operator; the two-mode combinations
$\phi_{2n}=\psi_{2n}\psi_0(0)-\psi_0\psi_{2n}(0)$ (\file{:744}); the vectors
$\E(\phi_n)$ and their Gram--Schmidt orthonormalisation $\varepsilon_n$; and the
near-null phenomenon itself, which is the paper's abstract.

The mechanism is also theirs, and it is stated in one line: \emph{the radical of
$\QW$ contains the range of $\E$} (\file{:181}, \file{:699}), because
$\F_\mu(\E f)$ carries a factor $\zeta(\tfrac12-is)$ and so vanishes at every
zero.  The prolate input is needed only to force the support into
$[\lambda^{-1},\lambda]$, and they are explicit that the construction ``only
uses the prolate vectors without any reference to $\QW_\lambda$''. \proved

\subsection{Index convention: the governing index is 4}
\label{sec:convention}

Connes--Consani write $\psi_{m,\lambda}(x)=\PS_{2m,0}(2\pi\lambda^2,x/\lambda)$,
so their $m$ is \emph{half} the prolate index.  Connes' survey writes
$\chi_k(\lambda)^2=\Lambda_{2k}(c)$ with $c=2\pi\lambda^2$, so his $\chi_2$ is
prolate index~$4$.  Both agree with the labelling of
Connes--Consani--Moscovici, whose $\psi^+_\ell:=h_{4\ell}-\frac{h_{4\ell}(0)}{h_0(0)}h_0$
is written in Hermite functions indexed in full.  Throughout this paper prolate
indices are \emph{full} indices, so the governing defect is $1-\Lambda_4$ and
Connes' $1-\chi_2=\tfrac12(1-\Lambda_4)+O((1-\Lambda_4)^2)$. \proved

\begin{remark}
The collision is real and has cost this project time.  It cannot be resolved
numerically: the identity
$h_\lambda(0)=\alpha h_{0,\lambda}(0)(\chi_4-\chi_0)/\chi_4$, which looks like a
discriminator, holds \emph{exactly} in arbitrary precision for every partner mode
$m\equiv0\bmod4$ and fails for every $m\equiv2\bmod4$.  It carries exactly one
bit --- the finite-Fourier phase --- and both candidate readings sit on the same
side of it.
\end{remark}

\subsection{The named open problem these objects are adjacent to}
\label{sec:adjacent}

Precisely: \emph{semilocal Weil positivity}, and within it, CCM's programme of
finding ``the semilocal analogue of the prolate operator''
(\cite{CCM2023} \file{mainc2m24fine.tex:197}--\file{:198}). \announced

The state of that programme, as of the principals' own most recent statements:
\begin{itemize}\itemsep2pt
\item Weil positivity is \emph{proved} at the archimedean place, for test
functions supported in $[2^{-1/2},2^{1/2}]$ \cite{CC2021}; Connes' 2026 survey
restates it with the same hypotheses (\file{rhready.tex:1198}--\file{:1200},
again in prose at \file{:1298}). \proved
\item The semilocal \emph{trace formula} is Connes 1999 and is not the gap
(\file{rhready.tex:1233}--\file{:1238}). \proved
\item The semilocal \emph{Sonin space} is stable under enlarging $S$
(\cite{CCM2023} Theorem~2). \proved
\item The candidate semilocal prolate operator
$\mathbf W_{\lambda,S}=(H+\tfrac12)^2+\lambda^2N_S$ is a definition; its Jacobi
coefficients for general $S$ are deferred to a forthcoming paper
(\file{:259}); a second candidate, via the Weil representation of the
metaplectic cover of $SL_2(\A_S)$, is deferred to the same (\file{:270}); and the
property that would make either usable --- negative eigenspace $=$ Sonin space
--- is called ``another constraint in the search'' (\file{:266}), not a theorem.
\announced
\item Neither deferred item had appeared as of 12 August 2026.  This is a
searched negative, not a proof of absence: author listings for all three
principals were enumerated, the survey's \LaTeX{} contains \emph{metaplectic}
zero times and \emph{Weil representation} zero times, and version histories were
checked.  Journal-only or personally-hosted work would not have been seen.
\end{itemize}

Two further named items, from the survey's own subsection \emph{Remaining steps}
(\file{rhready.tex:1174}): (i)~that the smallest eigenvalue of $\QW_\lambda$ is
simple with even eigenvector, and (ii)~that $k_\lambda:=\E(h_\lambda)$ is a good
enough approximation of $\theta_x$.  Neither is semilocal positivity, and the
survey does not list semilocal positivity at all.  This paper's \S\ref{sec:scale}
is adjacent to (ii) and does not deliver it: a Rayleigh-quotient bound would need
a spectral gap, and the near-radical is $\sim2\mu$-dimensional with all its
eigenvalues minuscule, so the gap is minuscule too.

% ==========================================================================
\section{The one proved sign, and why it does not travel}
\label{sec:onesign}
% ==========================================================================

\subsection{The archimedean chain, and where the sign enters}

\cite{CC2021} proves: for $g\in C_c^\infty(\R_+^*)$ with support in
$[2^{-1/2},2^{1/2}]$ and $\widehat g(i/2)=\widehat g(0)=0$,
\begin{equation}\label{eq:CC2021}
  W_\infty(g\star g^*)\;\ge\;\operatorname{Tr}\big(\vartheta(g)\,\Sonin\,\vartheta(g)^*\big),
\end{equation}
$\Sonin$ the orthogonal projection onto the Sonin space. \proved{} Seventy pages
go into it.  The sign, however, enters in one line, and it enters twice in the
same way: $P\widehat PP\ge0$, and $\Sonin=\Sonin^*=\Sonin^2$, so
$\operatorname{Tr}(\vartheta(g)\Sonin\vartheta(g)^*)=\|\Sonin\vartheta(g)^*\|_{HS}^2\ge0$.
\emph{An involution and a square.}  The authors say so themselves in the
abstract (\file{weil-compo.tex:86}).  Everything else --- the operator
$Q=-(\rho\partial_\rho)^2+\tfrac14$ implementing the vanishing conditions
without changing supports, Boas--Kac, the trace-remainder $\delta$, its jump at
$\rho=1$, the dihedral decomposition, the Toeplitz analysis of $\mathcal K_I$
--- is sign-preserving bookkeeping and remainder control.

It is worth saying plainly what supplies the sign there, because it is
\emph{not} the prolate functions.  They supply the labels of the
two-dimensional pieces in the dihedral decomposition, i.e.\ the bookkeeping that
makes the remainder computable.  \textbf{Prolate theory is the coordinate
system, not the mechanism.}

The endgame turns on three numbers, all specific to $I=[\tfrac12,2]$:
$\mathcal K_I$ has exactly one eigenvalue $>1$, and only just
($\lambda_{\max}=1.05158$); the next is $\lambda_2=0.686494$, giving a gap
$>0.2278$ on the complement; and the offending eigenvector satisfies
$\langle\zeta\mid\xi_0\rangle\approx0.94865$, i.e.\ the bad direction is $95\%$
aligned with the one admissible linear condition.  Change any of the three
materially and the proof stops. \proved{} \cite{CC2021}

So any semilocal analogue must produce all three: a finite-dimensional bad
subspace, a spectral gap on its complement, and an alignment between each bad
direction and an admissible condition --- \emph{uniformly in $\lambda$}, which
the archimedean case never had to face, since it was proved at a single
$\lambda$.  Meanwhile the dimension of the near-radical grows like $2\mu$.

\subsection{What survives at a finite place, and what does not}

Most of the machinery survives.  The twist $\Delta^{1/2}$ and the involution
$f\mapsto f^*$, Boas--Kac, and the operator $Q$ are statements about
$C_c^\infty(\R_+^*)$, the same algebra at every place; the trace-formula
representation is proved semilocally (Connes 1999); the radical of $\QW$
containing $\operatorname{ran}\E$ is a statement about the semilocal form
already; and the semilocal Sonin space exists and is stable.  Positivity of
$\operatorname{Tr}(\vartheta(g)\Sonin_S\vartheta(g)^*)$ survives trivially,
because it is positivity of a compression. \proved

Two things do not survive, and they are the two that carry the sign.

\begin{enumerate}\itemsep3pt
\item \textbf{There is no semilocal Slepian--Pollak.}  The archimedean argument
needs the angle operator to be diagonalised by an explicitly known family with a
known spectrum, and a differential operator commuting with it --- what Connes'
survey calls ``the miraculous existence, discovered by the Bell Labs group''
(\file{rhready.tex:1293}).  Semilocally $P^S_T$ and $\widehat P^S_W$ exist and
the trace formula holds; the spectrum of $P^S\widehat P^SP^S$ is \emph{not
known}, and producing the operator that would make it computable is the entire
content of the semilocal-prolate-operator programme. \announced

\item \textbf{There is no comparison term at all.}  No semilocal $\delta$, no
semilocal $\epsilon$, no compactness statement, no spectral analysis.  The 60
pages that were the work at the archimedean place have not been begun
semilocally.
\end{enumerate}

In one sentence: \emph{the archimedean proof buys its sign with a compression and
pays for it with a diagonalisation; semilocally the compression is free and the
diagonalisation does not exist.}

\subsection{The prime terms are exactly indefinite}
\label{sec:indefinite}

This is a small theorem and it excludes a mechanism, so it is stated as one.
Put $x=\log u$, $L=\log\mu$, $a=\log n$; then $V(n)$ of \eqref{eq:hilbert} is
the compression to $L^2([-\tfrac L2,\tfrac L2])$ of $n^{-1/2}(T_a+T_a^*)$, with
$T_a$ translation by $a$.

\begin{proposition}[\ours, elementary]\label{prop:indefinite}
$V(n)$ is unitarily equivalent to $-V(n)$.  Its spectrum is exactly symmetric
about $0$, at every $\lambda$, before any support or vanishing condition is
imposed, and
\[
  \|V(n)\|=2n^{-1/2}\cos\!\Big(\frac{\pi}{m+2}\Big),
  \qquad m=\lceil L/\log n\rceil-1 .
\]
\end{proposition}

\begin{proof}
Multiplication by $e^{i\pi x/a}$ is a unitary of $L^2([-\tfrac L2,\tfrac L2])$
and conjugates $T_a$ into $-T_a$, hence $V(n)$ into $-V(n)$.  The norm is the
Boas--Kac bound already used in \cite{CC2021} at \file{weil-compo.tex:842}.
\end{proof}

Verified numerically as well: spectrum symmetric to $10^{-15}$, and the closed
form matched at every row, on the full space and on the even sector that carries
$\sigma^+$ (\file{verify_semilocal_gap.py}, check 1; re-run on this pass).

Three consequences and one warning.

\begin{itemize}\itemsep2pt
\item \textbf{No compression of a projection can produce the prime terms.}  A
compression $\vartheta(g)\Pi\vartheta(g)^*$ is positive; $V(n)$ is symmetric
about zero.  The mechanism of \eqref{eq:CC2021} does not extend to them even in
principle.
\item \textbf{The prime terms do not weaken as $\lambda$ grows.}  $\|V(n)\|$
increases with $L$ towards $2n^{-1/2}$.  Adding places does not perturb a proved
inequality; it adds a full-strength indefinite term.
\item \textbf{The number of them grows} like $\pi(\mu)+O(\sqrt\mu)$.
\item \textbf{Warning, and it is the house rule applied to this very
proposition.}  ``Exactly indefinite'' is \emph{invariant} under
$W\mapsto-W$.  It says a sign cannot come from here; it says nothing about which
sign.  It is admissible as a negative result and must not be quoted as though it
explained the direction of an inequality.
\end{itemize}

\subsection{One finite place is degenerate; the difficulty is between places}
\label{sec:burnol}

It is natural to ask whether the missing spectral description exists at a single
finite place.  It does, it has since 1999, and its content is that the object
degenerates.

Burnol \cite{Burnol1999} takes \emph{Connes' own operator}
$\widetilde{P_\Lambda}P_\Lambda U(f)$ over a non-archimedean field, sets
$\Lambda=q^n$, $Q_n:=\widetilde{P_n}P_n$, and decomposes by the characters of
the unit group.  His Theorem~XI gives the kernel of $Q_n^\chi$ in closed form in
every case; what matters here is what the proof says about the two projections
separately.  For ramified $\chi$ they are projections onto coordinate subspaces
of the \emph{same} orthonormal basis.  For trivial $\chi$ with $2n\ge\delta$, in
Burnol's own words, ``$P_n$ cuts off $\mathcal M_n$, $\widetilde{P_n}$ cuts off
$\mathcal L_n$, \emph{so they commute}''.  The remaining case $0\le2n<\delta$ is
empty over $\Q_p$, where the different exponent $\delta$ is $0$. \proved

\begin{observation}[\ours; Burnol does not phrase it this way]
Over $\Q_p$ the two cutoff projections commute at every $\Lambda=p^n$.  A
commuting pair has $P\widehat PP=P\widehat P$, a projection, with spectrum in
$\{0,1\}$; Connes' angle $\Theta$ takes only the values $0$ and $\pi/2$.  There
are no non-trivial angles, no plunge region, and hence no prolate family to find.
\end{observation}

This does not transfer to the semilocal case, and the reason is definitional, so
it can be checked without doing any mathematics.  The semilocal projections are
defined ``using the module'' $\mathrm{Mod}_S(u)=\prod_{v\in S}|u|_v$
(\file{rhready.tex:1236}--\file{:1240}), so $P^S_T$ cuts off
$\{x:\prod_v|x_v|_v\le T\}$ --- one condition on a \emph{product}, not a
condition at each place.  Hence $P^S_T\neq\bigotimes_v P_{T_v}$, and CCM state
only an \emph{inclusion} $\eta_S(P_\lambda)\subset P^S_\lambda$
(\file{mainc2m24fine.tex:766}--\file{:767}).

\begin{quote}
\textbf{Consequence, stated flatly.}  Whatever difficulty the semilocal angle
operator has is located \emph{entirely in the coupling between places through
the product module}, not in any individual place.  Each finite place on its own
contributes a commuting pair and a one-dimensional Sonin space (CCM,
Prop.\ \texttt{soninppp}).
\end{quote}

It is worth recording that Connes 1999 substitutes a construction for a spectrum
at exactly this step: having given the archimedean plunge description in full,
he writes that for an arbitrary set of places ``we just use the same rule as in
the case of function fields'', i.e.\ the map $\psi\mapsto\psi\otimes1_R$
(\file{zeta.tex:2823}--\file{:2830}).  That map is CCM's $\eta_S$, which is
hilbertian and \emph{not} unitary.  This is the sharpest available statement of
the gap, and it is in Connes' own text.

Neither CCM nor the survey cites Burnol's p-adic scattering paper --- their
Burnol citations are the three archimedean/de~Branges papers --- which is why an
earlier absence-finding of ours missed it.  The one-place computation is outside
the corpus's citation graph.

% ==========================================================================
\section{Archimedean-only positivity fails, and it is not an artefact}
\label{sec:noplacewise}
% ==========================================================================

\subsection{Status: observed, not proved}
\label{sec:notheorem}

Connes--Consani report that the archimedean contribution alone ceases to be
positive past $L=\log2$, that positivity is restored by adding the prime $2$,
and that perturbing $2$ to $1.9999$ or $2.0005$ destroys it
(\file{Spectraltriples.tex:542}--\file{:600}).  Their own prose is unambiguous
about the register: ``The first fact we \emph{report from the numerical
computations}\ldots'' (\file{:543}); ``we \emph{test numerically}\ldots''
(\file{:177}); ``we \emph{report graphical evidence}'' (\file{:210}).  The verbs
are indicative --- they assert the fact --- but they never prove it, and there is
\emph{no theorem, proposition, lemma or corollary environment anywhere} in the
three sensitivity subsections, checked mechanically. \observed

Where else this was looked for, so the negative can be told from an unexamined
one: \cite{CC2021} never asserts sharpness of its Theorem~1, and two nearby
statements are easy to mistake for it --- the proved remark that $D\circ Q$ is
not negative on $[2^{-1},2]$ (a different object: the remainder, limit
$\approx+2.98699$), and the observation that $\lambda_{\max}(\mathcal K_I)<1$
``no longer holds true when the interval length gets closer to $\log2$'' (the
\emph{method} degrading, not positivity failing).  Its own reason for the
interval cuts the other way: the interval is chosen ``so that rational primes are
not involved'', i.e.\ for convenience.  \cite{CCM2023} does not mention the
figures.  Connes' 2026 survey restates Theorem~1 twice with the same hypotheses
and never mentions that the archimedean form goes negative past that range ---
if sharpness were known as a theorem, that is where it would appear.  The one
place it could be hiding is Yoshida \cite{Yoshida1992}, which we did not read.

\subsection{Their computation, rebuilt}

The third possibility --- that the failure is an artefact of truncation or of
double precision --- is excluded.  The even matrix $\sigma^+$ was rebuilt from
the explicit formula, with sign conventions taken from the formula rather than
inherited, and it matches Connes--Consani on all five numbers they report.

\begin{center}\small
\begin{tabular}{@{}p{0.42\textwidth}p{0.18\textwidth}p{0.32\textwidth}@{}}
\hline
Connes--Consani & source & ours \\
\hline
smallest eigenvalue at $L=\log2$, $\sim0.00133$ & \file{:551} &
$0.0013303$ at $N=80$; $0.0013296$ at $N=160$ \\
archimedean-alone crossing at $\mu=2.27$ & \file{:568} &
$\mu^*=2.2710$ at $N=160$, still decreasing in $N$ \\
$\infty+2$ at $\mu=3$: $<6\times10^{-8}$ & \file{:568} & $5.55\times10^{-8}$ \\
admissible $p$-window at $\mu=3$: width $<10^{-3}$ & \file{:574} &
$[1.9999995,\,2.00043]$, width $4.4\times10^{-4}$ \\
prime $3$ repairs past $\mu=3$ & \file{:579} &
at $\mu=3.5$: $\infty{+}2=-1.19\times10^{-2}$,
$\infty{+}2{+}3=+2.48\times10^{-10}$ \\
\hline
\end{tabular}
\end{center}

The second row carries the weight, because $2.27$ was not a number we aimed at:
the crossing was computed and then found written in their prose.  Their stated
window for it, $L\in[0.493,0.893]$, contains $L^*=0.82023$, in its ``upper
part'' as they say.

Why artefact is excluded, and it is not merely that the numbers agree.
\emph{Truncation cannot manufacture the failure}: restricting a quadratic form
to a subspace can only \emph{raise} its smallest eigenvalue, so a negative
truncated eigenvalue certifies a negative one for the full form.  Truncation can
hide a failure of positivity; it can never invent one.  The computed values
behave accordingly, decreasing monotonically in $N$.  \emph{Double precision is
not load-bearing for the sign change}: entries are $O(1)$, the spectrum at
$N=160$ runs $[1.3\times10^{-3},5.5]$, every number quadrature-converged to ten
digits --- four orders of magnitude away from the regime where this project's
arithmetic is unsound.  The rows that \emph{are} in that regime (the repaired
eigenvalues, $10^{-8}$ down to $10^{-12}$) are reported as reproductions, not
verifications, and nothing above leans on them: the load-bearing values are the
negative ones, of size $10^{-2}$ to $10^{-5}$.

\subsection{A witness, and why failure had to come eventually}
\label{sec:witness}

Two things the reproduction does not by itself give: a statement independent of
any matrix, and a reason.

\begin{proposition}[\ours]\label{prop:witness}
Let $f(x)=\cos(\pi x/L)+\kappa\cos(3\pi x/L)$ on $[-L/2,L/2]$ in the additive
coordinate $x=\log u$, continuous on $\R$ since both modes vanish at the
endpoints, with $\kappa$ fixed by $\int f(x)\cosh(x/2)\,dx=0$.  For real even $f$
that integral is $\widehat f(\pm i/2)$, so the boundary term of
\eqref{eq:hilbert} vanishes exactly and the archimedean value is $-W_\R$ alone.
It crosses zero at $\mu=3.5581$ and is negative at every larger $\mu$ tested;
at $\mu=4.0552$ it is $-0.646$ against $\|f\|^2=6.48$.
\end{proposition}

Evaluating it is three one-dimensional quadratures.  The threshold is weaker
than the true $2.2710$ --- it is one fixed vector, not the bottom eigenvector ---
but no truncation and no eigenvalue enters the argument, and the function can be
written down in a line.

\paragraph{The reason.}  On the Fourier side the archimedean part is
$\int|\widehat f(t)|^2\,\frac{2\vartheta'(t)}{2\pi}dt$ plus the boundary term,
and $\vartheta'$ is \emph{negative} on $|t|<6.2898$
($\vartheta'(0)=-2.68609170961\ldots$; first zero $6.2898360$).  So the Weil
density is negative in a neighbourhood of $t=0$, and the only thing stopping a
test function from living there is the uncertainty principle: a short support
forces $\widehat f$ to spread past the crossing, where $\vartheta'>0$.
Lengthening the support removes that obstruction.

This is a \emph{mechanism, not a proof} --- it says failure at large $L$ is
structurally expected and says nothing about where.  But it settles the shape of
the answer: the archimedean form is not a positive form that happens to be
verified on a small interval; it is an indefinite form whose positivity on a
small interval is an uncertainty-principle effect.  Turning it into a theorem ---
an explicit family with an explicit threshold --- looks routine and is not
attempted here.  \textbf{Gap G1} (\S\ref{sec:gaps}).

\begin{quote}
\textbf{Consequence for the shape of the missing theorem.}  Neither the
archimedean term beyond $\mu=2.2710$ nor any $-\Lambda(n)V(n)$ at any $\mu$ is a
positive form.  Only the sum is --- conjecturally, since that is RH.  So the
required theorem is about \emph{cancellation between places}, not about the
assembly of per-place positives.  That is a different and harder object than the
one that was proved.
\end{quote}

% ==========================================================================
\section{The archimedean deficit is bounded, and the balance is elsewhere}
\label{sec:deficit}
% ==========================================================================

Past the threshold, how does the archimedean deficit compare with what the
primes supply?  Write
\[
  \sigma^{\mathrm{arch}}=W_{0,2}-W_\R,\qquad
  D(\mu)=-\lambda_{\min}(\sigma^{\mathrm{arch}}),\qquad
  R(\mu)=-\sum_pW_p(v_{\rm arch}),
\]
with $v_{\rm arch}$ the minimiser of $\sigma^{\mathrm{arch}}$; so $D$ is the
deficit and $R$ is what the primes supply \emph{on the direction where the
archimedean side is worst}.  All computations in this section are in arbitrary
precision, which is not caution: the quantity that settles the competition is
$\sim10^{-48}$ at $\mu=11$ against matrix entries of order $1$, so double
precision cannot represent the answer, let alone resolve it.

\subsection{The bound}
\label{sec:bound}

\begin{theorem}[\ours, \proved{} given \cite{CC2023} Prop.\ \texttt{Hilbert}]
\label{thm:deficit}
For every $\mu$ and every unit $f$ in the even sector,
$\sigma^{\mathrm{arch}}(f,f)\ge2\vartheta'(0)\|f\|^2$; hence
\[
  D(\mu)\;<\;-2\vartheta'(0)\;=\;\log\pi-\psi(\tfrac14)\;=\;5.3721834192\ldots
\]
\end{theorem}

\begin{proof}
Three one-line facts applied to \eqref{eq:hilbert}.
(i) \emph{$\vartheta'$ attains its minimum at $t=0$.}
$\vartheta'(t)=\tfrac12\Re\psi(\tfrac14+\tfrac{it}2)-\tfrac12\log\pi$ and
$\frac{d}{dt}\Re\psi(\tfrac14+\tfrac{it}2)=-\tfrac12\Im\psi'(\tfrac14+\tfrac{it}2)$;
since $\psi'(z)=\sum_{k\ge0}(z+k)^{-2}$, at $z=\tfrac14+\tfrac{it}2$ each term
has imaginary part $-t(k+\tfrac14)/|z+k|^4<0$ for $t>0$.  So $\Im\psi'<0$,
$\vartheta'$ is strictly increasing on $t>0$, and
$\inf_t\vartheta'=\vartheta'(0)$.
(ii) \emph{On the even sector the boundary term is a square}: for real $f$ even
under $u\mapsto u^{-1}$, $\widehat f(\tfrac i2)=\widehat f(-\tfrac i2)=\int
F\cosh(x/2)dx$ is real, so the term is $2(\int F\cosh(x/2)dx)^2\ge0$.
(iii) \emph{Plancherel} in this normalisation:
$\int|\widehat f(t)|^2dt/2\pi=\|f\|^2$.
\end{proof}

What is not re-proved is Proposition \texttt{Hilbert} itself, which is proved in
\cite{CC2023}.  Its normalisation is the one place this could silently go wrong
by a factor of $2$, so it is \emph{checked and not assumed}: both sides of
\eqref{eq:hilbert} were evaluated on $f=\xi_1$, where $\widehat f$ is
closed-form, agreeing to $2\times10^{-8}$ at $\mu=2$ and $1\times10^{-8}$ at
$\mu=20$ (the residual is the $\log t/t^2$ tail, added from its asymptotic
form), with Plancherel confirmed at $0.999999998$.

That check is limited by its own test function: $\widehat{\xi_1}$ decays like
$t^{-1}$, so $|\widehat f|^2\vartheta'$ decays like $\log t/t^2$ and the tail has
to be supplied rather than measured.  The recheck of \S\ref{sec:recheck} removes
the difficulty by changing the test function to
$f=(1+\cos(a_1x))^2=\tfrac32+2\cos(a_1x)+\tfrac12\cos(a_2x)$, which vanishes to
second order at $\pm L/2$; its transform has both the $t^{-1}$ and the $t^{-3}$
coefficients cancelling identically, leaving
$\widehat f\sim12a_1^4\sin(tL/2)/t^5$, so the tail past a cutoff $T$ is
$O(\log T/T^9)$ and can simply be dropped.  At $\mu=8$, cutting off at $T=200$,
\[
  \Big(\int|\widehat f(t)|^2\tfrac{2\vartheta'(t)}{2\pi}\,dt\Big)
  \Big/\big(-W_\R(f,f)\big)\;=\;1.000000000000000007 ,
\]
and the residual is the tail rather than a constant: the same cutoff applied to
Plancherel --- which is exact in this normalisation --- leaves a residual of the
same size, raising the working precision from $40$ to $70$ digits moves neither,
and raising $T$ from $60$ to $200$ improves the ratio by a factor
$3.1\times10^4$ against the
$(200/60)^9=2.6\times10^4$ that $O(T^{-9})$ predicts.  \emph{There is no extra
factor of $2$, to $18$ digits.} \observed{} \ours{} (\S\ref{sec:recheck})

What remains taken on trust is the proposition itself, not its normalisation.
Their own proof of it says $\vartheta'$ is ``lower bounded'' and never extracts
the constant.

\subsection{It saturates}

\begin{center}\small
\begin{tabular}{@{}rrr@{}}
\hline
$\mu$ & $L$ & $D(\mu)$, $N=160$ \\
\hline
$3$ & $1.09861$ & $0.07397889$ \\
$5$ & $1.60944$ & $0.37270355$ \\
$11$ & $2.39790$ & $0.77466583$ \\
$20$ & $2.99573$ & $1.02682548$ \\
$10^{2}$ & $4.60517$ & $1.62332472$ \\
$10^{3}$ & $6.90776$ & $2.44603691$ \\
$10^{6}$ & $13.8155$ & $4.23599851$ \\
$10^{9}$ & $20.7233$ & $4.87322925$ \\
$10^{12}$ & $27.6310$ & $5.10422655$ \\
$10^{20}$ & $46.0517$ & $5.28349631$ \\
\hline
bound & & $5.37218341923\ldots$ \\
\hline
\end{tabular}
\end{center}

$D$ is increasing in $\mu$ --- as it must be, since the supports are nested and
the form is the same form, so $\lambda_{\min}$ can only fall.  Over the computed
range it grows like $\log\log\mu$; that is descriptive only and is not offered
as a law.  Whether $D(\mu)\to-2\vartheta'(0)$ exactly is open: \textbf{Gap G2}.

\begin{quote}
\textbf{Why this matters more than it looks.}  The natural reading of
``archimedean positivity fails past $\mu=2.271$, and the primes repair it'' is
that the thing to be repaired grows without limit and the repair must grow to
match.  It does not.  The archimedean side is negative by at most $5.3722$,
ever, on a unit vector.  \emph{The difficulty is not magnitude.}
\end{quote}

\subsection{On the archimedean bad direction the primes over-repair}

Since $s(\mu)\le\sigma^{\mathrm{arch}}(v_{\rm arch})-\sum_pW_p(v_{\rm arch})=R-D$
and $s>0$ under RH, positivity forces $R>D$; the question is by how much.

\begin{center}\small
\begin{tabular}{@{}rrrrr@{}}
\hline
$\mu$ & $D(\mu)$ & $R(\mu)$ & $R-D$ & $R/D$ \\
\hline
$3$ & $0.07393778$ & $0.12404899$ & $0.05011122$ & $1.6777$ \\
$4$ & $0.23972403$ & $0.26633090$ & $0.02660687$ & $1.1110$ \\
$5$ & $0.37262595$ & $0.41590755$ & $0.04328161$ & $1.1162$ \\
$6$ & $0.47506903$ & $0.48993904$ & $0.01487000$ & $1.0313$ \\
$8$ & $0.62435999$ & $0.64978693$ & $0.02542693$ & $1.0407$ \\
$10$ & $0.73104987$ & $0.74250899$ & $0.01145912$ & $1.0157$ \\
$12$ & $0.81337243$ & $0.83273410$ & $0.01936168$ & $1.0238$ \\
$16$ & $0.93617534$ & $0.94601094$ & $0.00983560$ & $1.0105$ \\
$20$ & $1.02672849$ & $1.04224568$ & $0.01551719$ & $1.0151$ \\
\hline
\end{tabular}
\end{center}

\noindent$R>D$ at every computed $\mu$ --- but read the last column rather than
the fourth: the \emph{relative} excess closes, from $1.678$ at $\mu=3$ to
$1.015$ at $\mu=20$, while the absolute slack oscillates in $0.010$--$0.050$
with maxima just after a new prime power enters. \observed

If the semilocal question were the question of these two columns it would be
answered in the affirmative, on a margin of $1.5\%$ at $\mu=20$ and falling.  It
is not that question.

\subsection{The balance is on a third direction, and its signs are reversed}
\label{sec:thirddirection}

The bottom eigenvector $v_\mu$ of the \emph{full} form is not $v_{\rm arch}$.
Splitting the Rayleigh quotient there, at $N=100$ and $120$ digits:

\begin{center}\small
\begin{tabular}{@{}rrrrrr@{}}
\hline
$\mu$ & $\sigma^{\mathrm{arch}}(v_\mu)$ & $-\sum_pW_p(v_\mu)$ & $s(\mu)$ &
$\log_{10}s$ & digits \\
\hline
$5$ & $+0.03958194$ & $-0.03958194$ & $1.005\times10^{-17}$ & $-16.998$ & $15.60$ \\
$6$ & $+0.04538231$ & $-0.04538231$ & $8.211\times10^{-23}$ & $-22.086$ & $20.74$ \\
$7$ & $+0.04959492$ & $-0.04959492$ & $7.731\times10^{-28}$ & $-27.112$ & $25.81$ \\
$8$ & $+0.05273312$ & $-0.05273312$ & $4.581\times10^{-33}$ & $-32.339$ & $31.06$ \\
$9$ & $+0.05520046$ & $-0.05520046$ & $3.135\times10^{-38}$ & $-37.504$ & $36.25$ \\
$10$ & $+0.05719477$ & $-0.05719477$ & $1.784\times10^{-43}$ & $-42.749$ & $41.51$ \\
$11$ & $+0.05882638$ & $-0.05882638$ & $1.226\times10^{-48}$ & $-47.912$ & $46.68$ \\
$12$ & $+0.06018733$ & $-0.06018733$ & $5.531\times10^{-54}$ & $-53.257$ & $52.04$ \\
\hline
\end{tabular}
\end{center}

\noindent(``digits'' is
$\log_{10}(|\sigma^{\mathrm{arch}}(v_\mu)|/s(\mu))$: the number of decimal
digits to which the two halves agree.)

\textbf{The signs are the other way round.}  On the direction that decides
positivity the archimedean contribution is \emph{positive} and the prime
contribution is \emph{negative}.  The phrase ``the primes repair the archimedean
deficit'' describes the previous subsection's direction and is \emph{false} on
this one, where the primes are the entire threat and the archimedean side is
what holds the form up.  Both statements are true of their own vector; what
follows is that the deficit-versus-repair columns are not the two sides of the
balance, and that the comfortable margin above buys nothing. \observed{}
(verified for $5\le\mu\le12$; the range is short --- \textbf{Gap G3}.)

\textbf{Both halves are small and grow slowly.}
$|\sigma^{\mathrm{arch}}(v_\mu)|$ is $0.0396$ at $\mu=5$ and $0.0602$ at
$\mu=12$; divided by $\log\mu$ it sits between $0.0242$ and $0.0255$ across the
range, so it is $\approx0.025\log\mu$.  Both halves are two orders below the
deficit $D$: the near-radical direction is not where the archimedean form is
bad, it is where it is nearly \emph{flat}.

\subsection{Every prime power is load-bearing}
\label{sec:loadbearing}

Connes--Consani show graphically, up to $\lambda^2\sim7$, that ignoring the
contribution of a prime power below $\lambda^2$ destroys positivity
(\file{:177}, figures at \file{:576}--\file{:600}).  Here that is a number.  At
$\mu=12$, $N=60$, $70$ digits, dropping one prime power at a time:

\begin{center}\small
\begin{tabular}{@{}lr@{}}
\hline
omitted & $\lambda_{\min}$ \\
\hline
none & $+9.19\times10^{-54}$ \\
$2$ & $-0.71226224$ \\
$3$ & $-0.83688017$ \\
$4$ & $-0.34656987$ \\
$5$ & $-0.71772943$ \\
$7$ & $-0.59699566$ \\
$8$ & $-0.17682307$ \\
$9$ & $-0.24427816$ \\
$11$ & $-0.10260697$ \\
\hline
\end{tabular}
\end{center}

\noindent(The ``none'' row is the $N=60$ upper bound; the $N=100$ value at
$\mu=12$ is $5.53\times10^{-54}$ --- consistent, since a looser truncation gives
a larger bound.)

Omit any single $p^m<\mu$ and the form goes negative by $0.10$ to $0.84$:
\emph{fifty-three orders of magnitude} above the residue that survives when all
are present.  There is no dominant prime, no ordering by size ($3$ costs more
than $2$; $11$, the largest, costs least but still $10^{52}$ times the residue),
and no tail that can be estimated crudely at any cutoff below $\mu$. \observed

% ==========================================================================
\section{The scale of \texorpdfstring{$s(\mu)$}{s(mu)}}
\label{sec:scale}
% ==========================================================================

\subsection{The index is forced, and it is 4}
\label{sec:index}

$\E$ is defined on $\mathcal S_0^{ev}$, i.e.\ $f(0)=\widehat f(0)=0$.  These two
conditions are not optional: they are what makes $\E f$ a test function at all.
For $f=\sum_mb_m\psi_{m,\lambda}$,
\[
  f(0)=\sum_mb_m\psi_m(0),\qquad
  \widehat f(0)=\int f=\sum_mb_m\chi_m\psi_m(0),
\]
the second because $\F_{e_\R}\psi_m=\chi_m\psi_m$ on $[-\lambda,\lambda]$
evaluated at $0$.  Two modes can satisfy both \emph{only if
$\chi_{m_1}=\chi_{m_2}$}; otherwise the $2\times2$ system is non-singular and
forces $b=0$.  Since $\chi_m\simeq(-1)^m$, that requires $m_1\equiv m_2\bmod2$
in Connes--Consani's labelling, i.e.\ \emph{prolate indices congruent mod 4}.

\begin{itemize}\itemsep1pt
\item $m=0$ alone: $\psi_0(0)\ne0$, inadmissible.
\item prolate indices $\{0,2\}$: $\chi\simeq+1$ and $-1$.  Excluded \emph{by the
phase, not by size}.
\item prolate indices $\{0,4\}$: both $\chi\simeq+1$.  Admissible, and minimal.
This is Connes--Consani's $\phi_2$.
\end{itemize}

\noindent So the governing defect is $1-\Lambda_4$, i.e.\ $1-\chi_2$ in Connes'
notation. \proved{} (\ours{} as an argument; the construction is
Connes--Consani's, and mode-4 selection by the finite-Fourier phase was found
earlier in this project in the Hermite rather than the prolate limit.)

\begin{remark}[a correction to Connes--Consani's construction, and it is free]
\label{rem:thirdmode}
Since $\chi_0=\chi_2$ only approximately, $\phi_2$ satisfies $\phi_2(0)=0$
exactly but
$\widehat{\phi_2}(0)=(\chi_0-\chi_2)\psi_0(0)\psi_2(0)\neq0$, with
$|\chi_0-\chi_2|\simeq1-\chi_2$.  That is not a rounding remark: by Poisson a
non-zero $\widehat f(0)$ puts an $x^{-1/2}$ tail on $\E f$ as $x\to0$, which is
not in $L^2(d^*u)$, and equivalently a pole in $\F_\mu(\E f)$ at $s=\pm\tfrac i2$
--- which is exactly what the two conditions are for.  It can be removed exactly
with a third mode: taking $\phi=b_0\psi_0+b_2\psi_2+b_4\psi_4$ (prolate indices
$0,4,8$, all with $\chi\simeq+1$) and imposing both conditions gives
$u_4/u_2\simeq-(1-\chi_2)/(1-\chi_4)$, a relative weight of $1.3\times10^{-8}$ at
$\mu=12$.  The vector is still essentially $\phi_2$ and the governing defect is
still $1-\Lambda_4$. \ours
\end{remark}

\subsection{Fuchs at general index, and Connes' constant}
\label{sec:fuchs}

Fuchs' asymptotic, in the form Connes cites it, is
$1-\Lambda_n(c)\sim4\sqrt\pi\,8^nc^{\,n+1/2}e^{-2c}/n!$.  Connes states
(\file{rhready.tex:1149})
\[
  1-\chi_2\;\sim\;\frac{2^{14}\sqrt2\,\pi^5}{3}\,\mu^{9/2}e^{-4\pi\mu},
\]
and with $1-\chi_2=\tfrac12(1-\Lambda_4)+O((1-\Lambda_4)^2)$ and $c=2\pi\mu$
this is \emph{identically} one half of Fuchs at $n=4$ --- verified to $100$
digits, difference $1.4\times10^{-101}$.  The $9/2$ is $n+\tfrac12$ at $n=4$.
\proved{} \theirs{Connes, from Fuchs}

Two things follow that matter for what any numerical agreement can be evidence
\emph{for}.

\begin{quote}
\textbf{The rate cannot discriminate the index.}  The exponential factor
$e^{-2c}$ in Fuchs \emph{does not depend on $n$}.  At $c=2\pi\mu$ every prolate
index predicts the same rate $2c/\mu=4\pi$.  A test that every candidate passes
is not a test.  What discriminates is the power $n+\tfrac12$ and the constant ---
and both are published. \ours{} as an observation
\end{quote}

\noindent Also, the asymptotic is not accurate enough to be used as the
comparison object.  Checked at $260$ digits against exact prolate eigenvalues,
the ratio asymptotic/exact at $n=4$ is $1.31219$ at $c=2\pi\cdot5$, $1.17936$ at
$c=2\pi\cdot8$, $1.11442$ at $c=2\pi\cdot12$, falling to $1.02897$ at $c=280$:
Fuchs \emph{over}-estimates by $1+O(1/c)$, and by $11$--$31\%$ over the whole
range where $s(\mu)$ is computed.  So the confrontation below uses exact
$1-\chi_2$.  (This is a check of the formula against the thing it approximates;
an earlier pass of ours had checked two \emph{formulas} against each other,
which is a weaker statement.) \ours

\subsection{The confrontation}
\label{sec:confront}

$s(\mu)$ recomputed at $N=100$, $120$ digits; $1-\Lambda_0$ and $1-\chi_2$ from
the prolate differential operator in the normalised Legendre basis, with the
apparatus validated three ways (Slepian's tabulated
$\Lambda_0(1)=0.57258$ recovered as $0.572581780638\ldots$; agreement to every
printed digit with an independent dense eigensolver; no movement across working
precisions $80$--$200$ digits and Legendre truncations $K=165,220$).

\begin{center}\small
\begin{tabular}{@{}rlllll@{}}
\hline
$\mu$ & $s(\mu)$, $N=100$ & $1-\Lambda_0$ & $s/(1-\Lambda_0)$ & $1-\chi_2$ &
$s/(1-\chi_2)$ \\
\hline
$5$ & $1.00502\times10^{-17}$ & $2.02073\times10^{-26}$ & $4.974\times10^{8}$ & $1.29839\times10^{-18}$ & $7.741$ \\
$6$ & $8.21103\times10^{-23}$ & $7.73821\times10^{-32}$ & $1.061\times10^{9}$ & $1.07935\times10^{-23}$ & $7.607$ \\
$7$ & $7.73054\times10^{-28}$ & $2.91979\times10^{-37}$ & $2.648\times10^{9}$ & $7.78995\times10^{-29}$ & $9.924$ \\
$8$ & $4.58127\times10^{-33}$ & $1.08993\times10^{-42}$ & $4.203\times10^{9}$ & $5.07899\times10^{-34}$ & $9.020$ \\
$9$ & $3.13549\times10^{-38}$ & $4.03552\times10^{-48}$ & $7.770\times10^{9}$ & $3.06721\times10^{-39}$ & $10.223$ \\
$10$ & $1.78439\times10^{-43}$ & $1.48462\times10^{-53}$ & $1.202\times10^{10}$ & $1.74467\times10^{-44}$ & $10.228$ \\
$11$ & $1.22587\times10^{-48}$ & $5.43359\times10^{-59}$ & $2.256\times10^{10}$ & $9.45817\times10^{-50}$ & $12.961$ \\
$12$ & $5.53055\times10^{-54}$ & $1.98020\times10^{-64}$ & $2.793\times10^{10}$ & $4.92908\times10^{-55}$ & $11.220$ \\
\hline
\end{tabular}
\end{center}

Against $1-\Lambda_0$ the ratio is $\sim10^9$ and \emph{grows by a factor $56$}
across the range --- almost all of that being the growth of
$(1-\chi_2)/(1-\Lambda_0)$ itself, which is $\mu^{4.3}$ here, since Fuchs' index
ratio contributes $\mu^4$.  The index-0 reading is not off by a constant; it is
off by the whole of the index-4 prefactor.  Against $1-\chi_2$ the ratio lies
between $7.6$ and $13.0$, with residual growth a factor $1.45$.

\emph{A factor of ten, nearly flat, against nine orders of magnitude with a
$\mu^4$ drift.}  That is the discrimination the rate could not make. \observed

One structure worth recording rather than smoothing: the ratio is largest at
$\mu=7,9,11$, each the last integer before a prime power enters the form.  $s$
is at its largest just before a new prime power arrives to repair it --- the
oscillation of $R-D$ seen from the near-radical direction, and a reminder that
$s(\mu)$ at integer $\mu$ samples a function with arithmetic structure at exactly
the scale of the sample spacing.

\subsection{The truncation bias, measured rather than asserted}
\label{sec:truncbias}

Truncation is variational and cuts one way, so every $s$ above is an
\emph{upper} bound, looser at larger $\mu$ --- which is the direction that
inflates $s/(1-\chi_2)$ at large $\mu$.  Computing $s(\mu)$ at $N=60,80,100,120$
and extrapolating geometrically in $N$ leaves ratios
$7.11,\ 9.65,\ 8.71,\ 9.92,\ 8.21,\ 9.21$ at $\mu=5,7,8,10,11,12$ --- a factor
$1.40$ spread with \emph{no monotone trend left in it}.  (At $\mu=6$ and $\mu=9$
the last two increments are not decreasing, so the geometric fit refuses and no
extrapolation is reported: that is the fit declining, not the sequence
diverging.)  At $N=100$ the ratio rose by $1.45$ across the range; after
extrapolation the rise is gone.

The extrapolation is a three-point model, not a theorem, but it can only move
the ratio \emph{down}, and it moves it down more at large $\mu$.  The claim made
here is therefore only the weak one: \emph{$\kappa=s/(1-\chi_2)$ grows no faster
than $\log\mu$}, which is also what \eqref{H1} would produce, since the mean
density of zeros near height $T$ is $\frac1{2\pi}\log\frac{T}{2\pi}$.  Either
way it is sub-polynomial and cannot touch an exponential rate.  Whether $\kappa$
is a constant of about $9$ is \textbf{Gap G4}.

\subsection{Their one published number}
\label{sec:theirnumber}

At $\mu=11$, Connes--Consani report $2.389\times10^{-48}$ (\file{:178}).
Convergence in $N$ here, at 100 digits: $2.389468\times10^{-48}$ at $N=50$;
$1.370549\times10^{-48}$ at $N=80$; $1.225866\times10^{-48}$ at $N=100$;
$1.120712\times10^{-48}$ at $N=120$.  \textbf{The exponent is settled at $-48$
and reproduces their number; the mantissa is not converged}, so this is a
reproduction and not a verification of $2.389$.

Note what $N=50$ does: it lands on their four printed digits.  Either that is
their truncation level or it is a coincidence; the paper never says which $N$
they used (\file{:210} says only $|n|,|m|\le N$), so this cannot be settled from
the paper and is recorded as unresolved.  On their own log-scale figures the
entire $N=50\to120$ drift is a change of $0.33$ against an ordinate of $-48$:
invisible, which is consistent with their remark that increasing $N$ ``does not
alter substantially the lower part of the spectrum'' without that remark being
evidence that the mantissa has converged.

\subsection{The identity, and the one inequality that is missing}
\label{sec:identity}

Everything in Connes--Consani's chain is exact except the last step, and the
last step is where a proof would live.  Here is what the chain actually gives.

Let $\phi$ be the normalised three-mode vector of Remark~\ref{rem:thirdmode}, so
that $\E\phi\in L^2(d^*u)$ and $\F_\mu(\E\phi)$ has no pole at $s=\pm\tfrac i2$.
Since $\operatorname{supp}\phi\subseteq[-\lambda,\lambda]$ we have
$\operatorname{supp}\E\phi\subseteq(0,\lambda]$, so with
\[
  g:=\E(\phi)\big|_{[\lambda^{-1},\lambda]},\qquad
  r:=\E(\phi)\big|_{(0,\lambda^{-1})},\qquad \E(\phi)=g+r,
\]
the decomposition is exhaustive: nothing spills above $\lambda$.

\begin{proposition}[\ours, from \cite{CC2023}]\label{prop:identity}
$\displaystyle \QW_\lambda(g,g)=\sum_{\frac12+is\in Z}\big|\F_\mu r(s)\big|^2 .$
\end{proposition}

\begin{proof}
$\F_\mu(\E\phi)$ vanishes on $Z$, so $\F_\mu g=-\F_\mu r$ on $Z$.
\end{proof}

This is an identity and it is the conceptual content: \textbf{$\QW_\lambda$
evaluated at the prolate vector is the Weil form of the part of $\E\phi$ that
falls outside the window.}  The near-radical is not ``nearly in the kernel'' by
accident; it is exactly as far from the kernel as the construction fails to fit
inside $[\lambda^{-1},\lambda]$.

\paragraph{The size of the spill.}  By Poisson, for $x<\lambda^{-1}$,
$\E\phi(x)=\E\widehat\phi(x^{-1})$, so with $y=x^{-1}>\lambda$,
$\|r\|^2=\int_\lambda^\infty|\sum_{n>0}\widehat\phi(ny)|^2dy$.  The $n=1$
diagonal term is half the out-of-band energy of $\phi$, and that energy splits
\emph{exactly} across the prolate modes, since
$\langle\widehat{\mathcal P_\lambda}\psi_n,\widehat{\mathcal P_\lambda}\psi_m\rangle=\Lambda_m\delta_{nm}$
makes the in-band parts orthogonal and hence the out-of-band parts too:
\[
  \big\|(1-\widehat{\mathcal P_\lambda})\phi\big\|^2
  =\sum_m|b_m|^2(1-\Lambda_{2m})=w\,(1-\Lambda_4)(1+o(1)),\qquad w\to\tfrac8{11}.
\]
So $\|r\|^2\asymp(1-\chi_2)$, linearly, and the corpus's central estimate
$\QW_\lambda(\E h_\lambda)\asymp1-\chi_4$ is the $n=1$ term of this line, with
the constant it hides being $\tfrac8{11}$ times an $O(1)$ mean value over the
zeros.  \emph{The $n\ge2$ terms and the cross terms are not free}: $\widehat\phi$
decays only like $y^{-1}$ off-band and the Cauchy--Schwarz sum is
logarithmically divergent before the oscillation of $\widehat\phi$ is used.  The
honest form is $\|r\|^2=C_1(\lambda)(1-\chi_2)$ with $C_1$ conjecturally
$O(\log\mu)$.  \textbf{Gap G5.}

\paragraph{The missing inequality.}  To convert this into a bound on $s$ one
needs two hypotheses, neither proved here:

\setcounter{hypothesis}{-1}
\begin{hypothesis}\label{H0}
$\|g\|^2$ is bounded below.  Since
$\|\E\phi\|^2=\frac1{2\pi}\int|\Gamma_\R(\tfrac12-is)\zeta(\tfrac12-is)P_\phi(s)|^2ds$,
this is a mean-value statement about $|\zeta(\tfrac12+it)|^2$ against an explicit
weight --- routine in shape, but $P_\phi$ depends on $\lambda$ and the bound must
be uniform.
\end{hypothesis}

\begin{hypothesis}\label{H1}
There is $\Theta(\lambda)$, subexponential in $\mu$, with
$\sum_{\frac12+is\in Z}|\F_\mu r(s)|^2\le\Theta(\lambda)\|r\|^2$ for the spill
$r$.
\end{hypothesis}

\eqref{H1} is a mean-value estimate for the zeta zeros as a sampling set, of
Plancherel--P\'olya type.  It is not automatic: $\QW_\lambda$ takes values in
$(-\infty,+\infty]$ and is unbounded above, so no uniform $\Theta$ exists on
$L^2$; \eqref{H1} must use the specific $r$, which is smooth on
$(0,\lambda^{-1})$ but has a jump at the endpoint, so $\F_\mu r$ decays only like
$|s|^{-1}$ and the sum converges only because the zero density is logarithmic.
The expected size of $\Theta$ is $O(\log\mu)$.

\begin{corollary}[conditional]\label{cor:upper}
Under \eqref{H0} and \eqref{H1},
\[
  s(\mu)\;\le\;\frac{\QW_\lambda(g,g)}{\|g\|^2}\;\le\;
  \frac{\Theta(\lambda)C_1(\lambda)}{\|g\|^2}\,(1-\chi_2(\lambda)),
  \qquad\text{hence}\qquad
  \limsup_{\mu\to\infty}\frac{\log s(\mu)}{\mu}\;\le\;-4\pi .
\]
\end{corollary}

That is a real half of the answer, and it is the half every number in this
project is consistent with, because every computed $s$ is itself an upper bound.
It also says that the residual truncation bias and the identification's
one-sidedness are the \emph{same} fact and not two.

% ==========================================================================
\section{The boundary: a lower bound of any shape implies RH}
\label{sec:boundary}
% ==========================================================================

A rate is two-sided.  The remaining half is that $s(\mu)$ does not decay
\emph{faster} than $e^{-4\pi\mu}$, i.e.\ a lower bound $s(\mu)\ge\kappa(\mu)>0$
for all $\mu$.

Connes states the equivalence himself (\file{rhready.tex:1145}): ``By the result
of Andr\'e Weil discussed in \S\texttt{sectweilpos}, the positivity of
$\QW_\lambda$ for all $\lambda>1$ is equivalent to RH''; Connes--Consani state
the same at \file{Spectraltriples.tex:169}. \proved{} Therefore:

\begin{theorem}[\ours{} as an observation, from Weil and \cite{Connes2026}]
\label{thm:boundary}
Any theorem of the form $\min(s^+(\mu),s^-(\mu))\ge\kappa(\mu)>0$ for all $\mu$
implies RH.  A theorem that additionally pins the rate is strictly stronger,
since RH gives positivity with no rate at all.
\end{theorem}

Two places one might hope to slip through, and neither works.

\begin{itemize}\itemsep3pt
\item \textbf{Restricting to large $\mu$ does not help.}  The supports are
nested, so $\QW_\Lambda\ge0$ implies $\QW_\lambda\ge0$ for every
$\lambda<\Lambda$: a lower bound valid for all $\mu\ge\mu_0$ already gives
positivity for \emph{all} $\lambda$.
\item \textbf{Dropping the sign does help --- and that is the tell.}  The
statement $|s(\mu)|=e^{-4\pi\mu+o(\mu)}$, with no claim that $s>0$, does not
formally imply RH.  But it is invariant under $\QW_\lambda\mapsto-\QW_\lambda$,
and is therefore compatible with the form being negative on the near-radical at
every $\mu$.  \emph{The only version of ``the rate is $4\pi$'' that says anything
about the sign of the Weil form is the version that contains RH.}
\end{itemize}

\begin{corollary}\label{cor:permanent}
The $4\pi$ rate cannot be made an unconditional theorem --- not for want of
technique; the statement contains RH.  And no computation can change that: a
wider grid or a converged truncation sharpens the fitted constant, but it cannot
supply a lower bound at a $\mu$ it has not computed, and the statement quantifies
over all $\mu$.
\end{corollary}

\begin{quote}
\textbf{The house rule predicted this.}  The mechanism of \S\ref{sec:objects}
--- the radical contains $\operatorname{ran}\E$, the near-radical is the failure
of a double support condition, the size is a concentration defect --- is entirely
invariant under $\QW_\lambda\mapsto-\QW_\lambda$.  A kernel is a kernel either
way; a concentration defect is a magnitude.  So the prolate story explains why
$|s|$ is minuscule and says nothing whatever about why $s>0$.  It is exactly that
sign-blindness which forces the lower bound to carry the whole arithmetic
content, and RH is the arithmetic content.  \emph{The standing test located the
obstruction before the analysis did.}
\end{quote}

For the record, the same two-sidedness has a numerical shadow that is worth
naming, because it is a check anyone can run: the three-parameter fit
$\log_{10}s(\mu)=-A\mu+B\log_{10}\mu+D$ over $\mu=5,\dots,12$ gives
$A=5.4635\pm0.052$, $B=5.322\pm0.96$, $D=6.589\pm0.44$, against the index-4
predictions $4\pi/\log10=5.45751$ ($0.12\sigma$), $9/2$ ($0.86\sigma$) and
$\log_{10}(2^{14}\sqrt2\pi^5/3)=6.373563$ ($0.49\sigma$) --- three parameters,
none free.  The index-0 reading predicts $B=\tfrac12$, which is $5.0\sigma$ away.
That is evidence for the \emph{identification}; by Corollary~\ref{cor:permanent}
it is not, and can never be, evidence for the theorem.

% ==========================================================================
\section{The obstruction, sharpened}
\label{sec:sharpening}
% ==========================================================================

This is the deliverable.  Combining \S\ref{sec:deficit} and \S\ref{sec:scale}: a
proof of $\QW_\lambda\ge0$ must establish that, on the near-radical direction,
\[
  \sigma^{\mathrm{arch}}(v_\mu)\;=\;\Big(\sum_pW_p\Big)(v_\mu)\;+\;s(\mu),
  \qquad
  \sigma^{\mathrm{arch}}(v_\mu)\approx0.025\log\mu,
  \qquad
  s(\mu)\approx\kappa\,(1-\chi_2(\lambda)),
\]
with $s>0$.  In words: \emph{two computable quantities of size
$\approx0.025\log\mu$, one archimedean and one a finite sum over $p^m<\mu$,
agree to about $5.2$ decimal digits per unit of $\mu$ over the computed range ---
at a rate consistent with $4\pi/\log10=5.4575$ --- and the sign of the
discrepancy is always the same one.}

Three consequences.  Each is a constraint on \emph{method}, not a restatement of
difficulty.

\begin{enumerate}\itemsep4pt
\item \textbf{No fixed-relative-error argument can work.}  Any bound on either
side accurate to a fixed relative error $\varepsilon$, or to a fixed power
$\mu^{-k}$, loses at
$\mu\gtrsim(k\log\mu+\log(1/\varepsilon))/(4\pi/\log10)$.  The required relative
accuracy is $e^{-4\pi\mu}/(0.025\log\mu)$: not a constant, and not polynomial.

\item \textbf{No prime can be dropped and no tail estimated.}  Each $p^m<\mu$
individually carries $0.10$--$0.84$, fifty-three orders above the residue
(\S\ref{sec:loadbearing}), and the cost is not ordered by the size of the prime.
The sum is not dominated by its head, and the standard move of controlling
$\sum_{p>P}$ crudely is unavailable at every $P<\mu$.

\item \textbf{The difficulty is not magnitude but cancellation, and
Theorem~\ref{thm:deficit} is what makes that precise.}  The deficit is bounded by
$5.3721834\ldots$ for all $\mu$ --- a fixed, small, explicit constant.  Nothing
on the archimedean side diverges.  What has to be proved is an
exact-to-$e^{-4\pi\mu}$ agreement between two bounded quantities, which is a
statement about the arithmetic of the primes and not about the size of anything.
\end{enumerate}

Read the other way, this is also an explanation of \emph{why} Connes--Consani's
near-radical construction is the right conceptual object: the near-radical is
exactly the subspace on which the two sides are forced into this agreement, and
its dimension $1+\nu(\mu)\sim2\mu$ is the number of independent such agreements
that $\mu$ demands at once.

And it is why the sharpening stops here rather than continuing.  A fourth
constraint would have to say something about the \emph{direction} of the
discrepancy; by \S\ref{sec:boundary} any such statement, of any shape, is RH.

% ==========================================================================
\section{The house rule, applied to this paper}
\label{sec:houserule}
% ==========================================================================

\begin{center}
\emph{Is any statement in this paper false for $-\QW_\lambda$?}
\end{center}

\noindent Applied claim by claim.  ``Sign-blind'' is not an accusation; it is a
classification, and the pattern in the last column is the paper's thesis.

\begin{center}\small
\begin{tabular}{@{}p{0.36\textwidth}p{0.38\textwidth}p{0.16\textwidth}@{}}
\hline
claim & under $\QW_\lambda\mapsto-\QW_\lambda$ & verdict \\
\hline
\S\ref{sec:objects}: the objects are Connes--Consani's; $\QW_\lambda$ is one
symbol; the index is 4
& statements about notation, about $\E$'s domain and about a phase; unchanged
& \textbf{sign-blind} \\
\S\ref{sec:onesign}: the archimedean sign is
$\operatorname{Tr}(\vartheta(g)\Sonin\vartheta(g)^*)\ge0$
& the theorem would read $-W_\infty(g\star g^*)\ge\operatorname{Tr}(\cdots)\ge0$,
contradicting the positivity Yoshida established by explicit computation
& \textbf{false for $-W$} \\
Prop.~\ref{prop:indefinite}: $V(n)\sim-V(n)$
& invariant by construction
& \textbf{sign-blind} \\
\S\ref{sec:burnol}: the projections commute over $\Q_p$
& a statement about two projections
& \textbf{sign-blind} \\
\S\ref{sec:noplacewise}: archimedean positivity fails past $\mu=2.2710$;
Prop.~\ref{prop:witness}'s witness has value $-0.646$
& becomes a claim about negativity being restored by the prime 2, and about a
value $+0.646$; both false of the same computation
& \textbf{false for $-W$} \\
Thm.~\ref{thm:deficit}: $\lambda_{\min}(\sigma^{\mathrm{arch}})\ge2\vartheta'(0)$
& becomes $\lambda_{\max}\le5.372$; measured $\lambda_{\max}=5.4855$ at $\mu=2$
and unbounded in $N$, since the form takes values in $(-\infty,+\infty]$
& \textbf{false for $-W$}, and not marginally \\
\S\ref{sec:deficit}: $R(\mu)>D(\mu)$
& the primes would have to reduce an archimedean excess and instead increase it
& \textbf{false for $-W$} \\
\S\ref{sec:thirddirection}: $\sigma^{\mathrm{arch}}(v_\mu)>0$,
$-\sum_pW_p(v_\mu)<0$, sum $=s>0$
& all three signs flip; the sum is $-s<0$
& \textbf{false for $-W$} \\
\S\ref{sec:loadbearing}: every prime power load-bearing
& ``the form goes negative'' becomes ``goes positive''
& \textbf{false for $-W$} \\
\S\ref{sec:confront}: $s/(1-\chi_2)\in[7,13]$
& $\lambda_{\min}(-\sigma^+)=-\lambda_{\max}(\sigma^+)$, large, negative, and
growing without bound in $N$
& \textbf{false for $-W$} \\
Prop.~\ref{prop:identity}: $\QW_\lambda(g,g)=\sum_Z|\F_\mu r|^2$
& both sides negate; the identity survives
& \textbf{sign-blind} \\
Cor.~\ref{cor:upper}: $\limsup\mu^{-1}\log s\le-4\pi$, and the rate $4\pi$
& magnitude only --- except for the left half of $0<s\le\cdots$, which fails
& \textbf{sign-blind} (the bound); \textbf{false for $-W$} (the positivity) \\
Thm.~\ref{thm:boundary}: a lower bound for all $\mu$ implies RH
& for $-W$ the hypothesis asserts $\sigma^+\prec0$, already false at $\mu=2$
where $\lambda_{\max}=5.4855$
& \textbf{false for $-W$} \\
\S\ref{sec:sharpening} constraint 1 (relative accuracy)
& a statement about digits
& \textbf{sign-blind} \\
\S\ref{sec:sharpening} constraints 2 and 3
& rest on \S\ref{sec:loadbearing} and Thm.~\ref{thm:deficit}, both sign-bearing
& \textbf{false for $-W$} \\
\S\ref{sec:recheck}: the numbers reproduce under a second implementation;
\S\ref{sec:erratum}: (\texttt{h02ev}) is a factor $2$ from the table at
\file{:444}
& statements about two computations agreeing and about two printed expressions;
a corroboration carries no direction of its own and inherits the verdict of
whatever row it corroborates
& \textbf{sign-blind} \\
\hline
\end{tabular}
\end{center}

Read the table.  \emph{Every sign-blind row that says anything about the form is
part of the prolate/magnitude machinery} --- the objects, the notation, the
index, the indefiniteness of the prime terms, the commuting pair at a finite
place, the identity, the rate.  (The last row says nothing about the form: it
records that two computations agree and that two printed expressions do not,
which no involution of $\QW_\lambda$ can touch.)
\emph{Every sign-bearing row is an inequality about $\lambda_{\min}$ of the
actual form}, or the one proved square at the archimedean place.  There is no row
in which the prolate mechanism produces a direction, and this is not for want of
looking: the classification was run on the claims as they stood.

The same test applied to the underlying note on the rate gave three sign-blind
claims out of six --- the index, the identity, and the rate itself --- with the
three sign-bearing ones each an inequality about $\lambda_{\min}$.  That is the
honest result and it is \emph{the argument}, not an embarrassment: it is what
turns ``we did not manage to prove positivity'' into ``this mechanism cannot
prove positivity, and here is the test that shows it''.

\paragraph{Two warnings to whoever quotes this paper.}
``$s(\mu)$ is governed by $1-\chi_2$'' is by itself invariant under
$W\mapsto-W$; it is sign-bearing only as ``$0<s(\mu)\le\cdots$'', and the left
half of that is RH.  Likewise ``the two halves agree to $5.46\mu$ digits'' is
sign-blind on its own, and is admissible only as part of ``\ldots and the
archimedean half is the positive one''.

% ==========================================================================
\section{Gaps, and open problems}
\label{sec:gaps}
% ==========================================================================

Recorded rather than filled; several were discovered while assembling this paper
and are not closed by it.

\begin{center}\small
\begin{tabular}{@{}lp{0.72\textwidth}@{}}
\hline
\textbf{G1} & Turn Prop.~\ref{prop:witness} into a theorem: an explicit family of
test functions and an explicit $\mu_0$ past which archimedean-only positivity
provably fails.  Looks routine.  It would upgrade the most-cited numerical claim
in this area from \observed{} to \proved. \\
\textbf{G2} & Is $D(\mu)\to-2\vartheta'(0)$ exactly, or to something strictly
smaller?  The saturation is numerical to $\mu=10^{20}$ and the gap is still
$0.089$.  The limit argument --- concentrate $\widehat f$ at $t=0$ while killing
the boundary term with an $O(e^{-L/4})$ correction near $\pm L/2$ --- is sketched
and not written out.  It would make Theorem~\ref{thm:deficit} two-sided. \\
\textbf{G3} & Is the sign reversal of \S\ref{sec:thirddirection} stable in $\mu$?
Verified only for $5\le\mu\le12$.  It is the load-bearing structural claim of
\S\ref{sec:deficit} and the range is short. \\
\textbf{G4} & Is $\kappa(\mu)=s/(1-\chi_2)$ a constant, and is it about $9$?
Extrapolated values are $7.1$--$9.9$ with no trend, but a three-point model over
eight points cannot separate a constant from $\log\mu$. \\
\textbf{G5} & Prove $\|r\|^2=C_1(\lambda)(1-\chi_2)$ with $C_1=O(\log\mu)$:
the $n\ge2$ terms and cross terms in \S\ref{sec:identity} need the oscillation of
$\widehat\phi$, since Cauchy--Schwarz alone diverges logarithmically. \\
\textbf{G6} & Prove \eqref{H1} --- the Plancherel--P\'olya-type mean-value bound
over the zeros for the spill $r$ --- and \eqref{H0}.  Together with G5 these give
$\limsup\mu^{-1}\log s\le-4\pi$ unconditionally.  \emph{This is the one genuinely
open analytic problem this paper produces}, and it is the whole remaining content
of the upper bound. \\
\textbf{G7} & The odd sector.  Connes--Consani's
$\phi_3=\psi_3\psi_1(0)-\psi_1\psi_3(0)$ is built from $\PS_{6,0}$ and
$\PS_{2,0}$, so the odd smallest eigenvalue should be governed by prolate index
\textbf{6} --- power $13/2$ at the same rate $4\pi$.  A sharp, cheap, falsifiable
prediction; not checked here. \\
\textbf{G8} & Evaluate $\QW_\lambda$ on $g$ directly and compare with $s(\mu)$.
Prop.~\ref{prop:identity} makes this the cleanest possible test of the whole
chain, replacing twelve figures with a ratio.  Not done: $\E\phi$ has jump
discontinuities at $u=\lambda/n$, so its coefficients decay like $j^{-1}$ and the
truncated Rayleigh quotient may be dominated by the tail rather than by the
$10^{-54}$ it is trying to measure.  Wants panel quadrature at the breakpoints. \\
\textbf{G9} & Does the corpus's commutator localisation have a semilocal
analogue, i.e.\ does $[\mathbf W_{\lambda,S},P^S]=0$ for either CCM candidate?
Answering it needs the deferred Jacobi coefficients. \\
\textbf{G10} & Not open, and recorded so it is not reopened: the matching
\emph{lower} bound on $s(\mu)$.  That is RH (Theorem~\ref{thm:boundary}). \\
\hline
\end{tabular}
\end{center}

% ==========================================================================
\section{Provenance, and how to check every claim}
\label{sec:provenance}
% ==========================================================================

\subsection{Sources read, and how}

All Connes / Connes--Consani / CCM / Burnol material was read as arXiv \LaTeX{}
source, downloaded from \texttt{arxiv.org/e-print/} on 12 August 2026; line
references of the form \file{Spectraltriples.tex:169} are lines of those files
and are stable for that snapshot.  Specifically:
\file{Spectraltriples.tex} (arXiv:2106.01715), \file{weil-compo.tex}
(arXiv:2006.13771), \file{mainc2m24fine.tex} (arXiv:2310.18423),
\file{rhready.tex} (arXiv:2602.04022), \file{zeta.tex} (arXiv:math/9811068),
\file{scatteringV2.tex} (arXiv:math/9901051).  Where only a statement was read
and not its proof, the underlying notes say so section by section.

\paragraph{Not read.}  Fuchs \cite{Fuchs1964} --- its statement is taken in the
form Connes cites and then checked against independent arbitrary-precision
computation of $\Lambda_0,\Lambda_2,\Lambda_4$ (\S\ref{sec:fuchs}); if the
general-$n$ constant were wrong, the $B$ and $D$ predictions of
\S\ref{sec:boundary} would move and the $A$ prediction would not.  Yoshida
\cite{Yoshida1992} --- not on arXiv and not reachable; it is the one place a
sharpness statement could hide (\S\ref{sec:notheorem}).  Slepian's tabulated
$\Lambda_0(1)=0.57258$ is quoted from the standard literature and used only as a
five-digit sanity check.

\subsection{Scripts}

Every numerical claim above has a companion script in the project's
\file{notes/} directory; all are self-contained, and the longer ones accept
\texttt{-{}-quick}.

\begin{center}\small
\begin{tabular}{@{}llp{0.42\textwidth}@{}}
\hline
script & needs & what it certifies here \\
\hline
\file{verify_semilocal_gap.py} & numpy &
Prop.~\ref{prop:indefinite} (symmetry to $10^{-15}$, Boas--Kac norm);
Fuchs vs Connes' constant to 12 decimals \\
\file{verify_arch_positivity.py} & numpy &
\S\ref{sec:noplacewise}: the five reproduced numbers, and
Prop.~\ref{prop:witness}'s witness ($\approx2.5$ min) \\
\file{verify_deficit_repair.py} & mpmath &
\S\ref{sec:deficit} entirely: the bound, the saturation table, $R/D$, the
Rayleigh split, the prime-omission table ($\approx20$ min; \texttt{-{}-quick}
$\approx4$) \\
\file{verify_prolate_rate.py} & mpmath &
\S\ref{sec:fuchs}--\S\ref{sec:theirnumber}: exact prolate defects, the
confrontation, the truncation extrapolation ($\approx40$ min;
\texttt{-{}-quick} $\approx5$) \\
\file{verify_index_convention.py} & mpmath &
\S\ref{sec:convention}: that the candidate discriminator carries one bit \\
\file{verify_independent_recheck.py} & nothing &
\S\ref{sec:recheck}: the five load-bearing numbers recomputed from the
definitions on the Python standard library alone; the (\texttt{thetaprime})
normalisation test of \S\ref{sec:bound}; the erratum of \S\ref{sec:erratum}
(\texttt{-{}-quick} drops to $N=60$ and loosens the bounds, which is the
variational point of \S\ref{sec:confront} showing up as a caveat) \\
\hline
\end{tabular}
\end{center}

\paragraph{What was re-run while assembling this paper, and what was not.}
Two scripts were re-run on 12 August 2026 (python 3.11, numpy 1.25.2, mpmath
1.3.0).  \file{verify_semilocal_gap.py} reproduced both of its checks:
Proposition~\ref{prop:indefinite}'s exact spectral symmetry and Boas--Kac norm at
every row, and the Fuchs-vs-Connes constant to twelve decimals.
\file{verify_prolate_rate.py --quick} reproduced check~0 (Slepian's
$\Lambda_0(1)=0.572581780638$; two independent eigensolvers agreeing to every
printed digit), check~1 (Connes' constant identified as $\tfrac12$ Fuchs at
$n=4$, and the Fuchs/exact ratios of \S\ref{sec:fuchs}), check~4 (the
three-parameter test) and check~5 (precision and truncation stability).  It did
\emph{not} reproduce the $s(\mu)$ column of \S\ref{sec:confront}: \texttt{-{}-quick}
runs $N\le80$, and its ratios ($8.7$ to $18.6$) sit above the $N=100$ ones, which
is the direction truncation must move them.

Everything else in \S\ref{sec:noplacewise}--\S\ref{sec:scale} is transcribed from
the notes named below, which each report their own script's output, and was
\emph{not} re-run on this pass.  That is the honest status of those numbers:
computed once, by the note that owns it, at the stated precision and truncation,
and reproducible by the command above.

One transcription error was found this way and is corrected here: the predicted
intercept $\log_{10}(2^{14}\sqrt2\pi^5/3)$ is $6.373563\ldots$, and one table in
\file{prolate-rate.md} gave $6.37347$.  The script had it right; the table did
not.  The $0.49\sigma$ and everything downstream are unchanged.  The note is
annotated in place.  The recheck of \S\ref{sec:recheck} recomputed that constant
from scratch and agrees with the correction rather than with the slip, which is
the only useful test a second pair of eyes can apply to a correction.

\subsection{The numerics, rechecked by a second implementation}
\label{sec:recheck}

Everything in \S\ref{sec:deficit} and \S\ref{sec:scale} was produced by one
implementation --- \file{verify_deficit_repair.py} and
\file{verify_prolate_rate.py}, written in one sitting by one author --- and
running the same code again is not a check of it.  Five of those numbers, the
load-bearing ones, have since been recomputed from the definitions in
\cite{CC2023} by a second implementation of ours,
\file{verify_independent_recheck.py}, written without opening either of the first
two: they were read only afterwards, once every number had been computed and
written down, which is the order the exercise depends on --- a reimplementation
written after reading the original inherits its slips silently and then agrees
with it.  \file{independent-recheck.md} reports the whole of it.  This subsection
gives the outcome, then what is independent in it, then what is not, the last of
these at the same length as the first because it is the part a reader cannot
reconstruct from the agreement.  Nothing here was re-run on this editing pass.

\paragraph{Outcome.}  All five load-bearing numbers reproduce, to every digit
printed here: the bound $|2\vartheta'(0)|=5.3721834\ldots$ by two routes and all
ten rows of the saturation table (\S\ref{sec:deficit}); all nine rows of $D$,
$R$, $R-D$, $R/D$, including the shape of the oscillation; all eight rows of the
reversed-sign split, signs and magnitudes (\S\ref{sec:thirddirection}); the
prime-omission table (\S\ref{sec:loadbearing}); the three fitted parameters,
their standard errors and all four $\sigma$ comparisons (\S\ref{sec:boundary});
and the prolate defects with both ratio columns (\S\ref{sec:confront}), giving
$s/(1-\chi_2)\in[7.6074,12.9609]$ and growth by a factor $56.16$ against
$1-\Lambda_0$, at $N=100$.  The $N$-convergence of \S\ref{sec:theirnumber}
reproduces to seven digits at all four truncations.  Two external anchors hold:
$\lambda_{\min}(\sigma^{\mathrm{arch}})=0.0013310$ at $\mu=2$ against
Connes--Consani's $\sim0.00133$ (\file{:551}), and Slepian's $\Lambda_0(1)$.

\paragraph{What is independent, stated precisely.}  The precise version is
checkable and ``independently verified'' is not.  The second implementation
shares with the first
\begin{itemize}\itemsep1pt
\item \textbf{no arithmetic library} --- Python's \texttt{decimal} (libmpdec,
decimal radix) against mpmath (binary radix over Python ints), so a rounding or
radix bug in one is not a bug in the other;
\item \textbf{no transcendental function} --- $\pi$, $\log$, $\exp$,
$\sin/\cos$, $\gamma$, $\psi$ and $\Re\psi$ of a complex argument are
implemented there from their series, where the original calls mpmath for every
one of them;
\item \textbf{no eigensolver} --- Householder tridiagonalisation and
Sturm-sequence bisection, which returns a \emph{bracket} whose endpoints are
separately checkable by an inertia count, against the original's iteration;
\item and, for the archimedean boundary term, \textbf{no method} --- it is
derived there in closed form as the rank-one matrix $W_{0,2}(n,m)=2v_nv_m$ with
$v_n=\int_{-L/2}^{L/2}\xi_n(x)e^{x/2}dx$, where the original computes that block
by quadrature.  That closed form is also what makes step (ii) of the proof of
Theorem~\ref{thm:deficit} --- the boundary term is a square --- visible in the
basis rather than asserted.
\end{itemize}
Two checks there have no counterpart in the original: the prolate normalisation
against the trace sum rule $\sum_n\Lambda_n(c)=2c/\pi$, to $58$ digits at $c=5$,
which pins the one factor a derivation of $\Lambda_n$ from the finite Fourier
transform could get wrong; and the convolution table at \file{:444} re-derived
twice, once symbolically and once by brute-force numerical convolution of the
basis functions themselves.

\paragraph{What is \emph{not} independent, and is therefore corroborated by none
of the above.}
\begin{itemize}\itemsep1pt
\item \textbf{The definitions.}  Both implementations read the same paper.
\emph{A shared misreading of \cite{CC2023} is invisible to this exercise}, and
under one every agreement above is worthless.  That is irreducible rather than a
defect of effort, and it is why the objects are tabulated by source line, here
and in the note, so that the reading itself can be disagreed with.  The only fix
is a third party who reads \cite{CC2023} without reading us; this is not a
substitute for one.
\item \textbf{Proposition \texttt{Hilbert} (\file{:289}) and eq.\
(\texttt{thetaprime}) (\file{:261}).}  Used as theorems by both implementations
and re-proved by neither.  The \emph{normalisation} of the second is now tested
numerically and holds to $18$ digits (\S\ref{sec:bound}); the proposition and
its proof are still taken from the paper by everyone here.
\item \textbf{The quadrature scheme.}  Composite Gauss--Legendre with
Newton-refined Legendre roots was chosen independently for the recheck and turned
out to be the same scheme the original uses.  That is a real limit on what the
agreement of the two matrices establishes, and \emph{arriving at the same scheme
independently is not the same thing as being independent of the scheme}.  What
does differ is the panel count, the node count and the grouping of the singular
integrand; varying panels and nodes moves $D(8)$ in none of $36$ decimals.
\item \textbf{The route to the prolate eigenvalues.}  The Legendre (Bouwkamp)
expansion and the identity $\Lambda_n=c\,d_0^2/(\pi\psi_n(0)^2)$ are the same in
both --- derived rather than copied, but the natural ones.  This is the weakest
independence claim of the exercise; the sum rule pins the normalisation and does
not pin the expansion.
\item \textbf{Fuchs' asymptotic} is quoted by both and derived by neither, as
\S\ref{sec:standards} already records.
\item \textbf{Not attempted at all}: the odd sector (\textbf{G7}); any re-proof
of Proposition \texttt{Hilbert}; and the mantissa of Connes--Consani's
$2.389\times10^{-48}$, which \S\ref{sec:theirnumber} leaves unverified and which
the recheck leaves unverified in exactly the same way.
\item \textbf{Not covered, of this paper's other numbers}: the geometric
extrapolation in $N$ of \S\ref{sec:truncbias} and the extrapolated ratios
$7.1$--$9.9$ that \S\ref{sec:whose} reports; the Fuchs-against-exact accuracy
ratios of \S\ref{sec:fuchs}; the one-bit discriminator of
\S\ref{sec:convention}; and everything in \S\ref{sec:onesign} and
\S\ref{sec:noplacewise}, which belong to a different note and a different
script.  Those stand where \S\ref{sec:provenance} leaves them --- computed once,
by the note that owns them.
\end{itemize}

\noindent And one thing the agreement is not, in this paper's own vocabulary
(\S\ref{sec:standards}): every $s(\mu)$ in both implementations is an
$N$-truncated \emph{upper} bound, so their agreement is agreement about a bound
and not about $s$.  It is a reproduction and not a verification --- the status
\S\ref{sec:confront} and Corollary~\ref{cor:upper} already give it, unchanged by
having been reproduced.

\paragraph{Three corrections, all cosmetic, and carried above.}
\S\ref{sec:deficit}'s $R(3)$ is $0.12404899$ and not $0.12404900$: the converged
value is $0.124048994835$, and $R-D$ and $R/D$ in that row were already right.
\S\ref{sec:thirddirection}'s ``digits'' entry at $\mu=8$ is $31.06$ and not
$31.02$: the row's own two entries give $\log_{10}(0.05273312)+32.33898=31.0611$,
and every other row of that column was right.  Both are corrected in the tables
above rather than left to the note.  The third is a definition and not a digit:
the residual rms of the three-parameter fit, quoted in \file{deficit-repair.md}
as $0.042$, is the $(n-p)$-normalised one, $0.04150$, the $/n$ one being
$0.03281$; this paper quotes only the fitted parameters and their standard
errors (\S\ref{sec:boundary}), which reproduce as
$A=5.463526\pm0.051591$, $B=5.322229\pm0.956951$, $D=6.588946\pm0.441716$ with
$A$ at $0.117\sigma$ from $4\pi/\log10$.  No claim in either note or in this
paper moves on any of the three.

\subsection{An erratum in \texorpdfstring{\cite{CC2023}}{CC2023}, and what does
not depend on it}
\label{sec:erratum}

This is here as a courtesy to a reader checking our arithmetic against
Connes--Consani's, and for nothing else.

Their Lemma \texttt{w02} gives closed forms for
$\widehat F(i/2)+\widehat F(-i/2)$ with $F(x)=\theta(\log x)$ and
$\theta=\xi_m\star\xi_n^*$.  The even-sector form as printed
(\file{Spectraltriples.tex:474}, eq.\ \texttt{h02ev}) is
\[
  \frac{8\,e^{-L/2}(e^{L/2}-1)^2L^3}{(L^2+16\pi^2m^2)(L^2+16\pi^2n^2)} .
\]
Carrying out the substitution the surrounding proof describes --- integrating
$(\theta(t)+\theta(-t))(e^{t/2}+e^{-t/2})$ over $[0,L]$ against the convolution
table at \file{:444}, from which the formula is derived four lines later ---
gives the same expression with $16$ in place of $8$.  Evaluated at $L=\log5$ for
$(n,m)=(1,2),(3,3),(4,7)$ the ratio of the two is $2.000000000000000$; the
odd-sector form at \file{:478} agrees with the same table at ratio
$1.000000000000000$.  Three routes give the $16$: term-by-term integration of the
\file{:444} even entry against $2\cosh(y/2)$; the rank-one form
$W_{0,2}(n,m)=2v_nv_m$ of \S\ref{sec:recheck}, which is the boundary term
$2\Re(\widehat f(\tfrac i2)\overline{\widehat f(-\tfrac i2)})$ of
\eqref{eq:hilbert} written in the basis, so that the factor is fixed by
Proposition \texttt{Hilbert} itself; and the $(0,0)$ entry, where
$\int_0^L(L-y)2\cosh(y/2)\,dy=4K$ exactly by parts with
$K=e^{L/2}-2+e^{-L/2}$. \observed{} \ours{} (elementary; the three routes are
written out in \file{independent-recheck.md})

\paragraph{The dependency, said out loud.}  \emph{Nothing in this paper depends
on the erratum being right.}  Neither of our implementations uses
(\texttt{h02ev}): both build $W_{0,2}$ from the table at \file{:444}, one by
quadrature and one from the closed form above, and the two agree.  What this
paper does depend on is our reading of that table --- if the table were being
read wrongly, the block $W_{0,2}$ and with it the whole of
\S\ref{sec:deficit} would move --- and that dependency is the shared-reading
exposure of \S\ref{sec:recheck}, not this one.  A reader who checks our
$\sigma^{\mathrm{arch}}$ against \file{:474} as printed will find it a factor
of~$2$ away, and that is the only reason this subsection exists.

\subsection{The notes this paper assembles}

In dependency order: \file{citation-audit.md} (positioning; $\QW_\lambda$ is one
symbol; the object-by-object correspondence); \file{index-convention.md} (the
index is 4, and the numerical route cannot settle it);
\file{semilocal-gap.md} (\S\ref{sec:onesign} and \S\ref{sec:noplacewise},
including the Burnol appendix and the archimedean-failure appendix);
\file{deficit-repair.md} (\S\ref{sec:deficit} and \S\ref{sec:sharpening});
\file{prolate-rate.md} (\S\ref{sec:scale} and \S\ref{sec:boundary});
\file{independent-recheck.md}, which depends on the last two and on nothing else
of ours (\S\ref{sec:recheck} and \S\ref{sec:erratum}).  Two earlier
notes --- \file{s3-reduction-audit.md} and \file{s3-sign-blindness.md} --- are
where the house rule of \S\ref{sec:houserule} was first stated and first used to
kill a chain of the project's own proposals; \file{signed-geometry-proposals.md}
is where it was sharpened to ``an involution and a square''.

\subsection{Corrections, recorded rather than absorbed}
\label{sec:corrections}

Four corrections of ours are folded into the text above.  They are listed here
because a reader weighing the tables should know which way the errors ran, and
because the pattern in them is itself a finding.

\begin{enumerate}\itemsep3pt
\item \textbf{A stated strategy read as a companion theorem.}  An early note
cited ``the archimedean place and its semilocal extension'' as though both were
proved.  Only the archimedean place is.  This is the origin of the status
vocabulary in \S\ref{sec:standards}.
\item \textbf{A check against one paper, written as a claim about the
literature.}  A row of the citation audit called the identification of the
controlling scale ``corpus-specific'' on the strength of arXiv:2106.01715.  It is
Connes', in a survey that predated the audit by six months and that the audit
itself named as the place to look.  A scoping defect, not a decayed finding.
\item \textbf{The wrong prolate index, from the nearer formula.}  The rate was
first read off the index-0 Slepian asymptotic, because that was the formula in
the note being continued --- while the correct index was already written into two
other notes of ours, two commits earlier.  Proximity beat provenance, and the
provenance was our own.
\item \textbf{A non-discriminating parameter mistaken for a test.}  The
$0.12\sigma$ agreement of the fitted rate with $4\pi$ was treated as evidence for
the index.  It is not: $e^{-2c}$ has no index dependence, so every candidate
passes.  \S\ref{sec:fuchs}.
\end{enumerate}

\noindent Three further corrections, all cosmetic and none of them of this kind,
came out of the recheck and are recorded with it in \S\ref{sec:recheck}.

\paragraph{The claim in this paper that would do the most damage if wrong} is the
ratio table of \S\ref{sec:confront}, because \S\ref{sec:index},
\S\ref{sec:identity} and \S\ref{sec:boundary} are all readings of it.  It has two
exposures: the prolate solver (guarded by an independent eigensolver and by
precision-stability checks), and the possibility that the near-flatness of
$s/(1-\chi_2)$ over $\mu=5..12$ is an accident of a short range --- eight integer
points spanning a factor $2.4$ in $\mu$, on a quantity that varies by $37$ orders
of magnitude across them.  The second is not guarded and cannot be at this cost.
It is, however, guarded against implementation error in the way described in
\S\ref{sec:recheck}, and against a shared misreading of \cite{CC2023} not at all.

The next most exposed was Theorem~\ref{thm:deficit}, whose only soft joint is the
normalisation of Connes--Consani's Proposition \texttt{Hilbert} --- a factor of
$2$ there doubles the constant.  \emph{That exposure is now closed}, and it is
recorded rather than deleted because a caveat that was real and was then
discharged is part of the evidence.  It was closed by test and not by argument:
eq.\ (\texttt{thetaprime}) is confirmed to $18$ digits on a test function whose
transform decays like $t^{-5}$, so that no asymptotic tail has to be supplied and
the residual is measurably the cutoff rather than a constant
(\S\ref{sec:bound}).  There is no extra factor of $2$.  What is still taken on
trust, and marked as such, is Proposition \texttt{Hilbert} itself: its
normalisation has been tested, its proof has not been re-proved by anyone here.

% ==========================================================================
\begin{thebibliography}{99}

\bibitem{Burnol1999}
J.-F. Burnol,
\emph{Scattering on the p-adic field and a trace formula},
Int. Math. Res. Not. \textbf{2000}, no.~2, 57--70; arXiv:math/9901051.

\bibitem{Connes1999}
A. Connes,
\emph{Trace formula in noncommutative geometry and the zeros of the Riemann zeta
function},
Selecta Math. (N.S.) \textbf{5} (1999), no.~1, 29--106; arXiv:math/9811068.

\bibitem{Connes2026}
A. Connes,
\emph{The Riemann Hypothesis: past, present and a letter through time},
arXiv:2602.04022 (Feb 2026).

\bibitem{CC2021}
A. Connes, C. Consani,
\emph{Weil positivity and trace formula, the archimedean place},
Selecta Math. (N.S.) \textbf{27} (2021), no.~4, Paper No.~77, 70 pp.;
arXiv:2006.13771.

\bibitem{CC2023}
A. Connes, C. Consani,
\emph{Spectral triples and $\zeta$-cycles},
Enseign. Math. \textbf{69} (2023), no.~1--2, 93--148; arXiv:2106.01715.

\bibitem{CM2022}
A. Connes, H. Moscovici,
\emph{The UV prolate spectrum matches the zeros of zeta},
Proc. Natl. Acad. Sci. USA \textbf{119} (2022), e2123174119; arXiv:2112.05500.

\bibitem{CCM2023}
A. Connes, C. Consani, H. Moscovici,
\emph{Zeta zeros and prolate wave operators: semilocal adelic operators},
Ann. Funct. Anal. \textbf{15} (2024), no.~4, Paper 87; arXiv:2310.18423.

\bibitem{CCM2025}
A. Connes, C. Consani, H. Moscovici,
\emph{Zeta spectral triples},
arXiv:2511.22755 (Nov 2025).

\bibitem{Fuchs1964}
W. H. J. Fuchs,
\emph{On the eigenvalues of an integral equation arising in the theory of
band-limited signals},
J. Math. Anal. Appl. \textbf{9} (1964), 317--330.
\emph{(Not read; used in the form cited by \cite{Connes2026} and checked
numerically --- \S\ref{sec:fuchs}.)}

\bibitem{Groskin2026}
A. Groskin,
\emph{A finite Guinand--Weil dictionary and archimedean tail order for the
truncated Weil quadratic form},
arXiv:2607.02828 (Jul 2026); and arXiv:2605.20224 (May 2026).

\bibitem{SlepianPollak1961}
D. Slepian, H. O. Pollak,
\emph{Prolate spheroidal wave functions, Fourier analysis and uncertainty --- I},
Bell System Tech. J. \textbf{40} (1961), 43--63.

\bibitem{Suzuki2026}
M. Suzuki,
\emph{Weil's quadratic form via the screw function},
arXiv:2606.09096 (Jun 2026).

\bibitem{Yoshida1992}
H. Yoshida,
\emph{On Hermitian forms attached to zeta functions},
Adv. Stud. Pure Math. \textbf{21} (1992), 281--325.
\emph{(Not read --- \S\ref{sec:notheorem}.)}

\end{thebibliography}

\end{document}
